\PassOptionsToPackage{table}{xcolor}
\documentclass[12pt,a4paper]{article}

% ====================================================================
% HỖ TRỢ TIẾNG VIỆT & FONT CHỮ SANG TRỌNG
% ====================================================================
\usepackage[utf8]{inputenc}
\usepackage[T5]{fontenc}
\usepackage[vietnamese]{babel}
\renewcommand{\refname}{Tài liệu tham khảo}
\usepackage{mathpazo} % Font Palatino chuyên nghiệp

% ====================================================================
% ĐỒ HỌA, HÌNH ẢNH & MÀU SẮC
% ====================================================================
\usepackage{amsmath, amsthm, amssymb}
\usepackage{graphicx}
\usepackage{tikz}
\usetikzlibrary{calc}
\usepackage{xcolor}
\definecolor{HCMUSBlue}{RGB}{0, 82, 155} % Màu xanh đặc trưng logo HCMUS
\definecolor{DeepBlue}{RGB}{10, 50, 110}
\definecolor{LightBlue}{RGB}{240, 245, 255}
\definecolor{GrayBg}{RGB}{245, 245, 247}
\definecolor{CommentGreen}{RGB}{34, 139, 34}
\definecolor{StringOrange}{RGB}{210, 105, 30}
\definecolor{KeywordBlue}{RGB}{0, 0, 205}

% ====================================================================
% HỘP ĐỊNH LÝ & GIẢ THIẾT NÂNG CAO (TCOLORBOX)
% ====================================================================
\usepackage[most]{tcolorbox}

% Hộp cho Định nghĩa / Giả thiết (Thanh dọc bên trái màu xanh thương hiệu)
\newtcolorbox{assumptionbox}[1]{
    enhanced,
    colback=HCMUSBlue!3,
    colframe=HCMUSBlue,
    leftrule=4pt,
    rightrule=0pt,
    toprule=0pt,
    bottomrule=0pt,
    arc=0pt,
    outer arc=0pt,
    title={\bfseries #1},
    coltitle=HCMUSBlue,
    detach title,
    before upper={\noindent\textbf{\color{HCMUSBlue}#1}\par\vspace{2mm}},
    boxsep=5pt,
    left=8pt,
    right=8pt,
    top=6pt,
    bottom=6pt
}

% Hộp cho Định lý / Mệnh đề toán học (Khung viền tròn bo góc nhẹ có tiêu đề nổi)
\newtcolorbox{theorembox}[1]{
    enhanced,
    colback=DeepBlue!2,
    colframe=DeepBlue,
    boxrule=1pt,
    arc=4pt,
    title={\bfseries #1},
    coltitle=white,
    attach boxed title to top left={yshift=-2mm,xshift=3mm},
    boxed title style={colback=DeepBlue,arc=2pt,boxrule=0pt},
    boxsep=6pt,
    left=8pt,
    right=8pt,
    top=8pt,
    bottom=8pt
}

% ====================================================================
% LAYOUT, SPACING & HỀ (GEOMETRY)
% ====================================================================
\usepackage{geometry}
\geometry{
    a4paper,
    left=30mm,
    right=20mm,
    top=30mm,
    bottom=30mm,
    headheight=15pt
}
\usepackage{setspace}
\setstretch{1.25} % Dãn dòng thoáng đãng, dễ đọc
\usepackage{indentfirst} % Thụt lề dòng đầu tiên của mỗi đoạn

% ====================================================================
% HEADER & FOOTER (FANCYHDR)
% ====================================================================
\usepackage{fancyhdr}
\pagestyle{fancy}
\fancyhf{}
\renewcommand{\headrulewidth}{0.6pt}
\renewcommand{\footrulewidth}{0.6pt}
\fancyhead[L]{\small \textit{Đồ án 2: Data Fitting và Phương Pháp OLS}}
\fancyhead[R]{\small \textit{Khoa CNTT - HCMUS}}
\fancyfoot[L]{\small \textit{Toán Ứng Dụng và Thống Kê}}
\fancyfoot[R]{\thepage}

% ====================================================================
% TIÊU ĐỀ PHÂN MỤC (TITLESEC)
% ====================================================================
\usepackage{titlesec}
\titleformat{\section}
  {\normalfont\large\bfseries\color{HCMUSBlue}}
  {\thesection}{1em}{}[{\titlerule[1.5pt]}]
\titleformat{\subsection}
  {\normalfont\normalsize\bfseries\color{DeepBlue}}
  {\color{HCMUSBlue}\ensuremath{\blacktriangleright}\ \thesubsection}{1em}{}

% ====================================================================
% BẢNG BIỂU, DANH SÁCH & HÌNH ẢNH
% ====================================================================
\usepackage{booktabs}
\usepackage{caption}
\usepackage{float}
\usepackage{enumitem}
\setlist[itemize]{noitemsep, topsep=4pt}
\usepackage{hyperref}
\hypersetup{
    colorlinks=true,
    linkcolor=HCMUSBlue,
    citecolor=HCMUSBlue,
    urlcolor=blue
}

% ====================================================================
% TRÌNH BÀY MÃ NGUỒN (LISTINGS)
% ====================================================================
\usepackage{listings}
\lstdefinestyle{pythonPremium}{
    language=Python,
    backgroundcolor=\color{GrayBg},
    basicstyle=\ttfamily\small,
    keywordstyle=\color{KeywordBlue}\bfseries,
    stringstyle=\color{StringOrange},
    commentstyle=\color{CommentGreen}\itshape,
    numbers=left,
    numbersep=8pt,
    numberstyle=\tiny\color{gray},
    frame=single,
    frameround=tttt,
    rulecolor=\color{HCMUSBlue!30},
    breaklines=true,
    breakatwhitespace=true,
    showstringspaces=false,
    captionpos=b,
    tabsize=4,
    keepspaces=true
}
\lstset{style=pythonPremium}

\begin{document}

% ====================================================================
% TRANG BÌA PREMIUM ĐẶC TRƯNG HCMUS
% ====================================================================
\begin{titlepage}
    \begin{tikzpicture}[remember picture, overlay]
        % Khung viền kép sang trọng
        \draw[line width=1.5pt, HCMUSBlue] 
            ([shift={(20mm,-20mm)}]current page.north west) 
            rectangle 
            ([shift={(-15mm,20mm)}]current page.south east);
        \draw[line width=0.5pt, HCMUSBlue!50] 
            ([shift={(22mm,-22mm)}]current page.north west) 
            rectangle 
            ([shift={(-17mm,22mm)}]current page.south east);
            
        % Họa tiết góc cổ điển
        \foreach \i in {north west, north east, south west, south east} {
            \node at (current page.\i) [shift={({\if\i\string north west 30mm\else\if\i\string north east -30mm\else\if\i\string south west 30mm\else-30mm\fi\fi\fi},{\if\i\string north west -30mm\else\if\i\string north east -30mm\else\if\i\string south west 30mm\else 30mm\fi\fi\fi})}] {
                \tikz[HCMUSBlue] \draw[line width=1pt] (0,0) (0,0.8) -- (0,0) -- (0.8,0);
            };
        }
    \end{tikzpicture}

    \begin{center}
        \vspace*{-1cm}
        {\large \textbf{ĐẠI HỌC QUỐC GIA THÀNH PHỐ HỒ CHÍ MINH}}\\[0.2cm]
        {\large \textbf{TRƯỜNG ĐẠI HỌC KHOA HỌC TỰ NHIÊN}}\\[0.2cm]
        {\large \color{HCMUSBlue} \textbf{KHOA CÔNG NGHỆ THÔNG TIN}}\\[1.0cm]
        \includegraphics[width=0.3\textwidth]{logo_hcmus.png}\\[1.0cm]
        
        \vfill
        
        \begin{tikzpicture}
            \node[inner sep=15pt, align=center, fill=HCMUSBlue!5, rounded corners=12pt] (titlebox) {
                \begin{minipage}{0.85\textwidth}
                    \centering
                    {\fontsize{20}{26}\selectfont \bfseries \color{HCMUSBlue} BÁO CÁO ĐỒ ÁN 2}\\[0.5cm]
                    {\fontsize{24}{30}\selectfont \bfseries \color{DeepBlue} Data Fitting và Phương Pháp OLS}\\[0.3cm]
                    {\fontsize{16}{20}\selectfont \bfseries \color{HCMUSBlue} Ứng dụng Dự báo chất lượng không khí PM2.5}\\[0.4cm]
                    \vspace{0.2cm}
                    {\large \itshape Môn học: Toán Ứng Dụng và Thống Kê}
                \end{minipage}
            };
            \draw[line width=2pt, HCMUSBlue] (titlebox.north west) -- (titlebox.north east);
            \draw[line width=2pt, HCMUSBlue] (titlebox.south west) -- (titlebox.south east);
        \end{tikzpicture}

        \vfill

        \begin{table}[H]
            \centering
            \renewcommand{\arraystretch}{1.2}
            \begin{tabular}{r l}
                \textbf{\large Nhóm thực hiện:}     & \textbf{\large Nhóm 6} \\
                \textbf{\large Sinh viên thực hiện:} & Nguyễn Quốc Đạt -- 24120282 \\
                                                     & Mai Văn Hiển -- 24120308 \\
                                                     & Nguyễn Đức Duy -- 24120294 \\
                                                     & Lê Quốc Hưng -- 24120184 \\
                                                     & Trần Nam Việt -- 24120155 \\
                \textbf{\large Lớp:}                 & 24CTT2 \\
                \textbf{\large Giảng viên hướng dẫn:} & ThS. Võ Nam Thục Đoan \\
                                                     & ThS. Lê Nhựt Nam \\
            \end{tabular}
        \end{table}

        \vfill
        \vspace{1cm}
        {\large \textbf{TP. Hồ Chí Minh, \the\year}}
        \vspace{1cm}
    \end{center}
\end{titlepage}

% ====================================================================
% MỤC LỤC
% ====================================================================
\newpage
\pagenumbering{roman}
\tableofcontents
\newpage

% ====================================================================
% NỘI DUNG BÁO CÁO
% ====================================================================
\pagenumbering{arabic}

\section{Giới thiệu chung}
Đồ án này tập trung vào nghiên cứu nền tảng lý thuyết và ứng dụng thực tiễn của hồi quy tuyến tính thông qua phương pháp bình phương tối thiểu (OLS) để giải quyết bài toán Data Fitting. Đồ án được chia làm hai phần chính:
\begin{enumerate}
    \item \textbf{Phần 1: Lý Thuyết Data Fitting và Minh Họa:} Nghiên cứu các giả thiết Gauss-Markov, xây dựng nghiệm OLS closed-form dưới dạng ma trận, chứng minh tính không chệch và tính tối ưu, cài đặt toàn bộ pipeline toán học từ đầu bằng Python thuần (không dùng thư viện số học nâng cao cho thuật toán chính) và minh họa thông qua mô phỏng Monte Carlo trên dữ liệu giả lập.
    \item \textbf{Phần 2: Ứng Dụng vào Thực Tế:} Áp dụng các kiến thức đã học vào phân tích và xây dựng mô hình dự báo chất lượng không khí (nồng độ bụi mịn PM2.5) dựa trên bộ dữ liệu khí tượng thực tế tại quận Shunyi, Bắc Kinh. Quy trình bao gồm khảo sát dữ liệu (EDA) chuyên sâu, thiết kế một pipeline tiền xử lý và loại bỏ đa cộng tuyến hoàn hảo bằng phân rã QR, xây dựng và so sánh 6 mô hình hồi quy (từ baseline tuyến tính đến các mô hình regularization và phi tuyến nâng cao), và tiến hành chẩn đoán phần dư thực tế.
\end{enumerate}

\clearpage

\section{Phần 1: Lý Thuyết Data Fitting và Minh Họa}

\subsection{Bài Toán Data Fitting}

\subsubsection{Phát biểu bài toán tổng quát}
Cho tập dữ liệu $\mathcal{D} = \{(x_i, y_i)\}_{i=1}^n$ với $x_i \in \mathbb{R}^p$ là vector đặc trưng và $y_i \in \mathbb{R}$ là biến mục tiêu. Bài toán data fitting là đi tìm hàm $f: \mathbb{R}^p \to \mathbb{R}$ trong một lớp hàm cho trước sao cho $f(x_i)$ xấp xỉ tốt nhất các giá trị thực tế $y_i$ theo một tiêu chí tối ưu cụ thể \cite{james}.

Trong mô hình hồi quy tuyến tính đa biến, ta giả thiết mối quan hệ có dạng:
\begin{equation}
    y_i = \beta_0 + \beta_1 x_{i1} + \beta_2 x_{i2} + \dots + \beta_p x_{ip} + \varepsilon_i = x_i^T \beta + \varepsilon_i
\end{equation}
Trong đó:
\begin{itemize}
    \item $\beta = (\beta_0, \beta_1, \dots, \beta_p)^T \in \mathbb{R}^{p+1}$ là vector hệ số hồi quy cần ước lượng (với $\beta_0$ là hệ số chặn).
    \item $x_i = (1, x_{i1}, \dots, x_{ip})^T \in \mathbb{R}^{p+1}$ là vector đặc trưng mở rộng của quan sát thứ $i$.
    \item $\varepsilon_i$ là sai số ngẫu nhiên (nhiễu).
\end{itemize}

Viết dưới dạng ma trận cho toàn bộ $n$ quan sát:
\begin{equation}
    y = X\beta + \varepsilon
\end{equation}
Với $y \in \mathbb{R}^n$, ma trận thiết kế $X \in \mathbb{R}^{n \times (p+1)}$ có cột đầu tiên toàn giá trị 1 để học hệ số chặn $\beta_0$, và $\varepsilon \in \mathbb{R}^n$.

\subsubsection{Các Giả Thiết Gauss-Markov}
Để phương pháp bình phương tối thiểu (OLS) có các tính chất thống kê tối ưu, ta cần thỏa mãn 5 giả thiết Gauss-Markov cơ bản sau \cite{hastie}:

\begin{assumptionbox}{Các Giả Thiết Gauss-Markov (GM1--GM5)}
\begin{itemize}
    \item \textbf{GM1 (Tính tuyến tính):} Biến phụ thuộc $y$ có mối quan hệ tuyến tính với các tham số $\beta$ và sai số ngẫu nhiên $\varepsilon$ theo công thức $y = X\beta + \varepsilon$.
    \item \textbf{GM2 (Không đa cộng tuyến hoàn hảo):} Ma trận thiết kế $X$ có hạng đầy đủ cột: $\operatorname{rank}(X) = p+1 < n$. Điều này đảm bảo các cột của $X$ độc lập tuyến tính hoàn toàn và ma trận đối xứng $X^T X$ khả nghịch.
    \item \textbf{GM3 (Tính ngoại sinh của sai số):} Kỳ vọng điều kiện của sai số bằng 0: $\mathbb{E}[\varepsilon \mid X] = \mathbf{0}$, tức là các biến độc lập không chứa thông tin để dự báo sai số.
    \item \textbf{GM4 (Đồng phương sai và không tự tương quan):} Ma trận hiệp phương sai của sai số có dạng:
    \begin{equation}
        \operatorname{Var}(\varepsilon \mid X) = \sigma^2 I_n
    \end{equation}
    Nghĩa là sai số tại mọi quan sát đều có phương sai không đổi $\sigma^2$ và sai số giữa hai quan sát bất kỳ độc lập với nhau (không có tự tương quan: $\operatorname{Cov}(\varepsilon_i, \varepsilon_j) = 0$ với mọi $i \neq j$).
    \item \textbf{GM5 (Nhiễu có phân phối chuẩn):} Sai số ngẫu nhiên phân phối chuẩn $\varepsilon \mid X \sim \mathcal{N}(\mathbf{0}, \sigma^2 I_n)$. Giả thiết này giúp thiết lập các khoảng tin cậy và thực hiện kiểm định giả thuyết thống kê chính xác ($t$-test, $F$-test).
\end{itemize}
\end{assumptionbox}

\subsection{Phương Pháp Bình Phương Tối Thiểu (OLS)}

\subsubsection{Hàm mất mát và nghiệm OLS}
Phương pháp OLS tìm kiếm vector tham số $\hat{\beta}$ sao cho tối thiểu hóa tổng bình phương phần dư (RSS) \cite{james}:
\begin{equation}
    RSS(\beta) = \|y - X\beta\|_2^2 = (y - X\beta)^T (y - X\beta) = y^T y - 2\beta^T X^T y + \beta^T X^T X \beta
\end{equation}

\begin{theorembox}{Định lý Nghiệm Bình Phương Tối Thiểu (OLS)}
Nếu ma trận thiết kế $X^T X$ khả nghịch, nghiệm cực tiểu duy nhất của bài toán OLS là:
\begin{equation}
    \hat{\beta}_{OLS} = (X^T X)^{-1} X^T y
\end{equation}
\end{theorembox}

\noindent\textit{Chứng minh:}
Lấy gradient của hàm mất mát $RSS(\beta)$ theo vector $\beta$:
\begin{equation}
    \nabla_{\beta} RSS(\beta) = \nabla_{\beta} (y^T y - 2\beta^T X^T y + \beta^T X^T X \beta) = -2X^T y + 2X^T X \beta
\end{equation}
Để tìm điểm cực tiểu, ta cho gradient bằng $\mathbf{0}$, dẫn đến hệ phương trình chuẩn (Normal Equations):
\begin{equation}
    X^T X \beta = X^T y
\end{equation}
Vì $X^T X$ khả nghịch (theo giả thiết GM2), nhân cả hai vế với ma trận nghịch đảo ta thu được nghiệm duy nhất:
\begin{equation}
    \hat{\beta}_{OLS} = (X^T X)^{-1} X^T y
\end{equation}
Để kiểm tra tính cực tiểu, ta tính ma trận Hessian: $\nabla^2_{\beta} RSS(\beta) = 2X^T X$. Do $X^T X$ là ma trận đối xứng xác định dương (vì các cột của $X$ độc lập tuyến tính), ma trận Hessian xác định dương, chứng tỏ nghiệm tìm được là điểm cực tiểu toàn cục duy nhất.

\noindent\textbf{Cài đặt toán học thực tế:}
Trong thực tế tính toán số học, việc tính trực tiếp ma trận nghịch đảo $(X^T X)^{-1}$ là cực kỳ không ổn định và dễ gây tích tụ sai số làm tròn. Trong mã nguồn \texttt{ols\_implementation.py}, hệ phương trình chuẩn $(X^T X) \beta = X^T y$ được giải trực tiếp bằng phương pháp \textbf{Khử Gauss chọn trục xoay một phần (Gaussian Elimination with Partial Pivoting)}. Thuật toán tìm phần tử chốt lớn nhất trên cột hiện tại để hoán vị hàng trước khi triệt tiêu phần tử dưới đường chéo, tránh chia cho các số cực nhỏ gây mất ổn định số học. Toàn bộ hàm giải hệ \texttt{manual\_solve(A, b)} được cài đặt thủ công bằng cấu trúc danh sách Python thuần (list-of-lists) mà không sử dụng bất kỳ thư viện bên ngoài nào.

\subsubsection{Ma Trận Chiếu và Hat Matrix}
Giá trị dự đoán của mô hình OLS được viết dưới dạng:
\begin{equation}
    \hat{y} = X\hat{\beta}_{OLS} = X(X^T X)^{-1} X^T y = H y
\end{equation}
Trong đó, ma trận $H = X(X^T X)^{-1} X^T \in \mathbb{R}^{n \times n}$ được gọi là \textbf{ma trận Hat} vì nó trực quan hóa việc đặt "mũ" ($\hat{}$) lên đầu vector quan sát $y$ \cite{strang}.

\noindent\textbf{Các tính chất quan trọng của $H$:}
\begin{itemize}
    \item \textbf{Tính đối xứng:} $H^T = (X(X^T X)^{-1} X^T)^T = X((X^T X)^{-1})^T X^T = X(X^T X)^{-1} X^T = H$.
    \item \textbf{Tính lũy đẳng (Idempotent):} $H^2 = H \cdot H = X(X^T X)^{-1} X^T X (X^T X)^{-1} X^T = X(X^T X)^{-1} X^T = H$. Ý nghĩa hình học: chiếu hai lần liên tiếp lên một không gian cột của $X$ cũng giống như chiếu một lần.
    \item \textbf{Giá trị riêng:} Chỉ nhận giá trị 0 hoặc 1.
    \item \textbf{Hạng ma trận:} $\operatorname{rank}(H) = \operatorname{Tr}(H) = p+1$.
    \item \textbf{Leverage:} Phần tử nằm trên đường chéo $h_{ii} = x_i^T (X^T X)^{-1} x_i$ biểu thị giá trị leverage của quan sát $y_i$ lên chính giá trị dự đoán $\hat{y}_i$.
\end{itemize}

\begin{theorembox}{Định lý Gauss-Markov}
Dưới các giả thiết Gauss-Markov GM1--GM4, ước lượng bình phương tối thiểu $\hat{\beta}_{OLS}$ là ước lượng tuyến tính không chệch tốt nhất, nghĩa là có phương sai nhỏ nhất trong tất cả các ước lượng tuyến tính không chệch khác:
\begin{enumerate}
    \item \textbf{Tính tuyến tính:} $\hat{\beta}_{OLS}$ là hàm tuyến tính của $y$: $\hat{\beta}_{OLS} = (X^T X)^{-1} X^T y = W y$.
    \item \textbf{Tính không chệch:} Kỳ vọng của ước lượng bằng giá trị thực:
    \begin{equation}
        \mathbb{E}[\hat{\beta}_{OLS} \mid X] = \beta
    \end{equation}
    \item \textbf{Tính tối ưu (phương sai nhỏ nhất):} Với mọi ước lượng tuyến tính không chệch khác $\tilde{\beta} = C y$, ma trận hiệp phương sai của nó luôn lớn hơn hoặc bằng ma trận hiệp phương sai của $\hat{\beta}_{OLS}$:
    \begin{equation}
        \operatorname{Var}(\tilde{\beta} \mid X) \succeq \operatorname{Var}(\hat{\beta}_{OLS} \mid X) = \sigma^2 (X^T X)^{-1}
    \end{equation}
\end{enumerate}
\end{theorembox}

\subsubsection{Ước Lượng Phương Sai Nhiễu}
Phương sai nhiễu $\sigma^2$ là một tham số chưa biết. Ước lượng không chệch của $\sigma^2$ được tính bằng tổng bình phương phần dư chia cho bậc tự do thực tế của sai số \cite{james}:
\begin{equation}
    \hat{\sigma}^2 = \frac{RSS}{n - p - 1} = \frac{\|y - X\hat{\beta}\|_2^2}{n - p - 1}
\end{equation}
Việc chia cho $n - p - 1$ nhằm bù đắp lượng thông tin đã mất khi ước lượng $p+1$ hệ số của vector $\hat{\beta}$ trước đó, đảm bảo kỳ vọng toán học $\mathbb{E}[\hat{\sigma}^2] = \sigma^2$.

\subsection{Đánh Giá Mô Hình}

\subsubsection{Các Chỉ Số Đo Lường Sai Số Dự Báo}
Hiệu năng dự báo định lượng của mô hình trên tập Test được đánh giá thông qua các thước đo sai số cơ bản:
\begin{itemize}
    \item \textbf{Sai số bình phương trung bình (MSE):} $MSE = \frac{1}{n} RSS$.
    \item \textbf{Căn sai số bình phương trung bình (RMSE):} $RMSE = \sqrt{MSE}$.
    \item \textbf{Sai số tuyệt đối trung bình (MAE):} $MAE = \frac{1}{n} \sum |y_i - \hat{y}_i|$.
\end{itemize}

\subsubsection{Hệ số xác định $R^2$ và $R^2$ hiệu chỉnh}
Hệ số xác định $R^2$ đo lường tỷ lệ phần trăm sự biến thiên của biến phụ thuộc $y$ được giải thích bởi các biến độc lập trong mô hình \cite{james}:
\begin{equation}
    R^2 = 1 - \frac{RSS}{TSS}
\end{equation}
Trong đó $TSS$ là tổng biến thiên toàn bộ của $y$ quanh giá trị trung bình. Do $R^2$ luôn có xu hướng tăng khi thêm biến, ta sử dụng \textbf{$R^2$ hiệu chỉnh} để phạt độ phức tạp của mô hình:
\begin{equation}
    \bar{R}^2 = 1 - \frac{n - 1}{n - p - 1} (1 - R^2)
\end{equation}

\subsubsection{Kiểm Định Giả Thuyết Thống Kê}
Dưới giả thiết nhiễu chuẩn GM5, ta tiến hành các kiểm định giả thuyết ý nghĩa thống kê:
\begin{itemize}
    \item \textbf{Kiểm định t:} Kiểm tra ý nghĩa của từng hệ số hồi quy riêng lẻ. Giả thuyết $H_0: \beta_j = 0$. Thống kê kiểm định t là $t_j = \hat{\beta}_j / (\hat{\sigma} \sqrt{[(X^T X)^{-1}]_{jj}})$. Nếu p-value hai phía của thống kê $t_j$ nhỏ hơn mức ý nghĩa $\alpha = 0.05$, ta bác bỏ $H_0$, chứng tỏ đặc trưng thứ $j$ đóng góp ý nghĩa vào mô hình.
    \item \textbf{Kiểm định F:} Kiểm tra ý nghĩa đồng thời của toàn bộ mô hình hồi quy. Giả thuyết $H_0: \beta_1 = \dots = \beta_p = 0$. Thống kê kiểm định F là:
    \begin{equation}
        F = \frac{(TSS - RSS) / p}{RSS / (n - p - 1)} \sim F_{p, n-p-1}
    \end{equation}
    Nếu p-value tương ứng cực kỳ nhỏ, ta bác bỏ $H_0$, xác nhận mô hình hồi quy có ý nghĩa giải thích tốt hơn mô hình chỉ có hệ số chặn.
\end{itemize}

\subsection{Các Vấn Đề Nâng Cao trong Hồi Quy}

\subsubsection{Đa cộng tuyến và chỉ số VIF}
Đa cộng tuyến xảy ra khi các biến độc lập có sự tương quan tuyến tính mạnh với nhau, khiến ma trận $X^T X$ gần suy biến và phương sai các hệ số $\operatorname{Var}(\hat{\beta}) = \sigma^2 (X^T X)^{-1}$ phình to. Hệ số phóng đại phương sai (VIF) dùng để phát hiện đa cộng tuyến cho biến thứ $j$ \cite{james}:
\begin{equation}
    VIF_j = \frac{1}{1 - R_j^2}
\end{equation}
Với $R_j^2$ là hệ số xác định khi hồi quy biến $X_j$ theo các biến còn lại. Nếu $VIF_j > 10$, đặc trưng $X_j$ bị đa cộng tuyến nghiêm trọng.

\subsubsection{Hồi Quy Điều Chuẩn (Ridge và Lasso)}
Để giải quyết đa cộng tuyến và chống overfitting, ta thêm thành phần phạt vào hàm mất mát \cite{hastie}:
\begin{itemize}
    \item \textbf{Hồi quy Ridge (chuẩn $L_2$):} Phạt tổng bình phương hệ số:
    \begin{equation}
        \hat{\beta}_{ridge} = \arg\min_{\beta} \left( \|y - X\beta\|_2^2 + \lambda \|\beta\|_2^2 \right) = (X^T X + \lambda I)^{-1} X^T y
    \end{equation}
    \item \textbf{Hồi quy Lasso (chuẩn $L_1$):} Phạt tổng trị tuyệt đối hệ số nhằm tạo nghiệm thưa thớt:
    \begin{equation}
        \hat{\beta}_{lasso} = \arg\min_{\beta} \left( \|y - X\beta\|_2^2 + \lambda \|\beta\|_1 \right)
    \end{equation}
\end{itemize}

\subsubsection{Chẩn đoán phần dư}
Phân tích phần dư kiểm chứng giả thiết Gauss-Markov \cite{james}:
\begin{enumerate}
    \item \textbf{Residuals vs Fitted:} Kiểm tra tính tuyến tính và phương sai không đổi.
    \item \textbf{Normal Q-Q Plot:} Kiểm tra giả thiết sai số phân phối chuẩn.
    \item \textbf{Scale-Location Plot:} Phát hiện xu hướng thay đổi phương sai.
    \item \textbf{Cook's Distance Plot:} Phát hiện các điểm dữ liệu dị biệt có ảnh hưởng lớn:
    \begin{equation}
        D_i = \frac{\hat{\varepsilon}_{i,std}^2 \cdot h_{ii}}{p \cdot (1 - h_{ii})}
    \end{equation}
\end{enumerate}

\subsubsection{Cross-Validation}
Phương pháp xác thực chéo k-Fold dùng để đánh giá độ tổng quát hóa của mô hình và tối ưu siêu tham số. Điểm số CV được tính bằng trung bình các sai số trên $k$ tập kiểm tra:
\begin{equation}
    CV_{(k)} = \frac{1}{k} \sum_{i=1}^k MSE_i
\end{equation}

\subsection{Minh họa Định lý Gauss-Markov bằng mô phỏng Monte Carlo}
Để kiểm chứng thực nghiệm định lý Gauss-Markov khẳng định OLS là ước lượng tuyến tính không chệch tốt nhất, ta tiến hành mô phỏng Monte Carlo $N = 1000$ lần trên dữ liệu giả lập được thiết kế như sau:
\begin{enumerate}
    \item Sinh ma trận thiết kế $X$ kích thước $100 \times 3$ cố định.
    \item Thiết lập vector tham số thực $\beta_{true} = [2, -1.5, 4]^T$.
    \item Tại mỗi vòng lặp $i \in \{1, \dots, N\}$, sinh vector nhiễu chuẩn $\varepsilon^{(i)} \sim \mathcal{N}(0, 0.5^2 I_{100})$, tính biến mục tiêu $y^{(i)} = X\beta_{true} + \varepsilon^{(i)}$.
    \item Tính ước lượng OLS: $\hat{\beta}_{OLS}^{(i)} = (X^T X)^{-1} X^T y^{(i)}$.
    \item Thiết lập một ước lượng tuyến tính không chệch thay thế (Alternative Estimator) có dạng $\tilde{\beta} = (W_{OLS} + C) y$ với $C$ là ma trận thỏa mãn $CX = \mathbf{0}$. Ma trận $C$ được khởi tạo bằng cách chiếu một ma trận nhiễu ngẫu nhiên lên không gian trực giao của $X$.
    \item \textbf{Phân tích kết quả thực nghiệm:}
    \begin{itemize}
        \item \textbf{Tính không chệch:} Kỳ vọng thực nghiệm $\frac{1}{N} \sum_{i=1}^N \hat{\beta}_{OLS}^{(i)}$ hội tụ chính xác về giá trị thực $\beta_{true} = [2, -1.5, 4]^T$.
        \item \textbf{Tính tối ưu:} Phương sai của ước lượng OLS luôn nhỏ hơn phương sai của ước lượng thay thế tại tất cả các hệ số ($\operatorname{Var}(\hat{\beta}_{OLS}) < \operatorname{Var}(\tilde{\beta})$). Trực quan hóa histogram cho thấy phân phối của OLS hẹp hơn và có đỉnh cao hơn rõ rệt so với ước lượng thay thế, chứng minh OLS có phương sai bé nhất.
    \end{itemize}
\end{enumerate}

\clearpage

\section{Phần 2: Ứng Dụng Data Fitting vào Dữ Liệu Thực Tế}

\subsection{Khảo sát \& Mô tả dữ liệu}

\subsubsection{Nguồn gốc và đặc điểm bộ dữ liệu}
Bộ dữ liệu được trích xuất từ bộ dữ liệu chất lượng không khí PRSA công bố công khai trên kho lưu trữ \href{https://archive.ics.uci.edu/dataset/501/beijing+multi+site+air+quality+data}{UCI Machine Learning Repository} \cite{ucidata}. Dữ liệu ghi nhận chất lượng không khí hàng giờ liên tục trong 4 năm (từ 01/03/2013 đến 28/02/2017) tại \textbf{trạm Shunyi, Bắc Kinh.}

\begin{table}[H]
\centering
\caption{Thông tin tổng quan tập dữ liệu Shunyi}
\rowcolors{2}{HCMUSBlue!5}{white}
\begin{tabular}{ll}
\toprule
\textbf{Thuộc tính} & \textbf{Giá trị} \\
\midrule
Thời gian bắt đầu & 01/03/2013 (00:00) \\
Thời gian kết thúc & 28/02/2017 (23:00) \\
Tần suất đo lường & Hàng giờ \\
Tổng số lượng quan sát ($n$) & 35,064 dòng \\
Tổng số lượng cột ($p$) & 18 cột \\
\bottomrule
\end{tabular}
\end{table}

\subsubsection{Mô tả chi tiết các biến số}
\begin{table}[H]
\centering
\small
\caption{Mô tả các đặc trưng trong tệp dữ liệu gốc}
\rowcolors{2}{HCMUSBlue!5}{white}
\resizebox{\textwidth}{!}{
\begin{tabular}{llll}
\toprule
\textbf{Tên cột} & \textbf{Kiểu dữ liệu} & \textbf{Loại biến} & \textbf{Mô tả chi tiết} \\
\midrule
\texttt{No} & int64 & Định lượng chỉ số & Số thứ tự của bản ghi (không dùng huấn luyện) \\
\texttt{year} & int64 & Định tính thứ bậc & Năm thực hiện đo lường (2013--2017) \\
\texttt{month} & int64 & Định tính thứ bậc & Tháng trong năm (1--12) \\
\texttt{day} & int64 & Định tính thứ bậc & Ngày trong tháng (1--31) \\
\texttt{hour} & int64 & Định tính thứ bậc & Giờ trong ngày (0--23) \\
\texttt{PM2.5} & float64 & Định lượng liên tục & Nồng độ bụi mịn PM2.5 ($\mu$g/m$^3$) - \textbf{Biến mục tiêu} \\
\texttt{PM10} & float64 & Định lượng liên tục & Nồng độ bụi thô PM10 ($\mu$g/m$^3$) \\
\texttt{SO2} & float64 & Định lượng liên tục & Nồng độ khí Sulfur Dioxide ($\mu$g/m$^3$) \\
\texttt{NO2} & float64 & Định lượng liên tục & Nồng độ khí Nitrogen Dioxide ($\mu$g/m$^3$) \\
\texttt{CO} & float64 & Định lượng liên tục & Nồng độ khí Carbon Monoxide ($\mu$g/m$^3$) \\
\texttt{O3} & float64 & Định lượng liên tục & Nồng độ khí Ozone ($\mu$g/m$^3$) \\
\texttt{TEMP} & float64 & Định lượng liên tục & Nhiệt độ không khí ($^\circ$C) \\
\texttt{PRES} & float64 & Định lượng liên tục & Áp suất khí quyển (hPa) \\
\texttt{DEWP} & float64 & Định lượng liên tục & Nhiệt độ điểm sương ($^\circ$C) \\
\texttt{RAIN} & float64 & Định lượng liên tục & Lượng mưa tích lũy (mm) \\
\texttt{wd} & object (str) & Định tính định danh & Hướng gió bề mặt (16 hướng phân loại) \\
\texttt{WSPM} & float64 & Định lượng liên tục & Tốc độ gió trung bình (m/s) \\
\texttt{station} & object (str) & Định tính định danh & Tên trạm đo = "Shunyi" (hằng số) \\
\bottomrule
\end{tabular}
}
\end{table}

\subsection{Khảo Sát Dữ Liệu Chuyên Sâu (EDA)}

\subsubsection{Cấu trúc dữ liệu và Thống kê mô tả}
Tiến hành phân tích cấu trúc dữ liệu và tính toán các chỉ số thống kê mô tả nâng cao bao gồm hệ số biến thiên ($cv\% = std/mean \times 100$), độ lệch skewness và độ nhọn kurtosis:
\begin{figure}[H]
    \centering
    \includegraphics[width=0.9\textwidth]{cell1.1.png}
    \caption{Phân tích kiểu dữ liệu và số lượng giá trị duy nhất trong tập dữ liệu}
\end{figure}

\begin{figure}[H]
    \centering
    \includegraphics[width=0.9\textwidth]{cell2.png}
    \caption{Thống kê mô tả chi tiết các đặc trưng định lượng}
\end{figure}

\subsubsection{Phân phối các biến số}
Để có cái nhìn trực quan về hình dáng phân phối, đồ thị histogram được vẽ cho toàn bộ các đặc trưng định lượng:
\begin{figure}[H]
    \centering
    \includegraphics[width=0.95\textwidth]{cell3.png}
    \caption{Đồ thị Histogram phân phối của các đặc trưng định lượng}
\end{figure}

\noindent\textbf{Nhận xét phân phối từ trực quan hóa:}
\begin{itemize}
    \item \textbf{Nhóm chất ô nhiễm và bụi mịn (\texttt{PM2.5}, \texttt{PM10}, \texttt{SO2}, \texttt{NO2}, \texttt{CO}, \texttt{O3}):} Phân phối lệch phải cực kỳ mạnh mẽ ($skewness > 1.5$). Trị số trung vị luôn nhỏ hơn nhiều so với giá trị trung bình. Điều này cho thấy phần lớn thời gian nồng độ chất ô nhiễm được kiểm soát ở mức thấp, nhưng sự hiện diện của phần đuôi kéo dài phản ánh các đợt bão bụi hoặc ô nhiễm không khí cực đoan xảy ra cục bộ theo mùa.
    \item \textbf{Nhóm nhiệt động học mùa vụ (\texttt{TEMP}, \texttt{DEWP}):} Phân phối có dạng hai đỉnh (bimodal). Đây là minh chứng vật lý rõ nét cho sự phân hóa thời tiết sâu sắc theo mùa tại Bắc Kinh (mùa hè nóng ẩm tạo đỉnh bên phải và mùa đông khô lạnh tạo đỉnh bên trái).
    \item \textbf{Nhóm khí tượng ngẫu nhiên (\texttt{PRES}, \texttt{WSPM}, \texttt{RAIN}):} Áp suất khí quyển (\texttt{PRES}) phân phối chuẩn rất rõ nét. Tốc độ gió (\texttt{WSPM}) lệch phải nhẹ cho thấy đa số là ngày lặng gió hoặc gió nhẹ (1--2 m/s). Lượng mưa (\texttt{RAIN}) tập trung tuyệt đối ở mốc 0.0 do thời gian tạnh ráo chiếm đại đa số.
\end{itemize}

\subsubsection{Phát hiện Outliers}
Vẽ biểu đồ Boxplot nhằm xác định sự phân bố của các giá trị dị biệt trong dữ liệu:
\begin{figure}[H]
    \centering
    \includegraphics[width=0.9\textwidth]{cell4.1.png}
    \caption{Đồ thị Boxplot xác định giá trị dị biệt của các đặc trưng chất ô nhiễm}
\end{figure}

Kết quả phân tích cho thấy các đặc trưng ô nhiễm môi trường (\texttt{SO2}: 8.9\%, \texttt{CO}: 6.0\%, \texttt{PM2.5}: 2.7\%) và các hiện tượng khí tượng ngẫu nhiên (\texttt{WSPM}: 5.7\%, \texttt{RAIN}: 4.7\%) chứa tỷ lệ outliers rất cao. Đây hoàn toàn là các biến động thực tế của môi trường (các đợt nghịch nhiệt tích tụ bụi hoặc bão lớn), do đó ta chọn phương pháp \textbf{Winsorization} (co các điểm dị biệt về ranh giới $Q_3 + 1.5 \times IQR$) thay vì xóa bỏ để giữ lại tín hiệu của các đợt cực trị.

\subsubsection{Ma trận tương quan Heatmap}
\begin{figure}[H]
    \centering
    \includegraphics[width=0.85\textwidth]{cell5.png}
    \caption{Ma trận hệ số tương quan tuyến tính Pearson}
\end{figure}

\noindent\textbf{Các tương quan vật lý nổi bật:}
\begin{itemize}
    \item Tương quan cực mạnh giữa \texttt{PM2.5} và \texttt{PM10} ($r \approx 0.88$) và \texttt{CO} ($r \approx 0.75$), xác nhận chúng có chung nguồn gốc phát thải (khói bụi công nghiệp, giao thông và sưởi ấm đốt than mùa đông).
    \item Tương quan thuận rất cao giữa nhiệt độ \texttt{TEMP} và nhiệt độ điểm sương \texttt{DEWP} ($r \approx 0.87$).
    \item Tương quan âm giữa \texttt{PM2.5} và tốc độ gió \texttt{WSPM} ($r \approx -0.27$), thể hiện hiệu ứng pha loãng vật lý (gió lớn giúp phát tán bụi mịn nhanh chóng).
\end{itemize}

\subsubsection{Biến động theo thời gian, chu kỳ mùa và giờ}
\begin{figure}[H]
    \centering
    \includegraphics[width=0.95\textwidth]{cell6.png}
    \caption{Biến động nồng độ bụi mịn PM2.5 trung bình hàng tháng giai đoạn 2013-2017}
    \label{fig:cell6}
\end{figure}

\begin{figure}[H]
    \centering
    \includegraphics[width=0.9\textwidth]{cell7.png}
    \caption{Phân phối nồng độ PM2.5 theo Mùa (trái) và theo Giờ trong ngày (phải)}
    \label{fig:cell7}
\end{figure}

Dựa vào các biểu đồ phân tích trực quan hóa chuỗi thời gian, phân phối theo mùa và biến động theo giờ của nồng độ bụi mịn \texttt{PM2.5} tại trạm Shunyi giai đoạn 2013--2017, ta có thể rút ra các nhận định và đánh giá chi tiết như sau:

\noindent\textbf{1. Biến động theo thời gian và xu hướng tổng thể (Hình \ref{fig:cell6}):}
\begin{itemize}
    \item \textbf{Tính chu kỳ mùa vụ (Seasonality) thể hiện rất mạnh mẽ:} Nồng độ \texttt{PM2.5} luôn tạo thành các đỉnh nhọn vọt lên rất cao vào chu kỳ cuối năm cũ - đầu năm mới (thường rơi vào các tháng 12, 1, 2). Các đỉnh này đạt mức cực kỳ nghiêm trọng, dao động từ $100$ đến trên $150\,\mu\text{g/m}^3$. Ngược lại, nồng độ có xu hướng giảm xuống mức thấp nhất vào mùa hè (tháng 6, 7, 8), tạo thành vùng trũng dao động quanh mức $40 - 60\,\mu\text{g/m}^3$.
    \item \textbf{Tình trạng ô nhiễm nghiêm trọng so với tiêu chuẩn WHO:} Đường biểu diễn nồng độ trung bình tháng \texttt{PM2.5} \textbf{chưa từng có một lần nào} giảm xuống dưới ngưỡng an toàn dài hạn của Tổ chức Y tế Thế giới (WHO AQG: $35\,\mu\text{g/m}^3$). Đồng thời, trong giai đoạn thu đông, nồng độ liên tục vượt xa ngưỡng cảnh báo 24 giờ của WHO (WHO Guideline 24h: $75\,\mu\text{g/m}^3$) với tần suất dày đặc.
    \item \textbf{Xu hướng tổng thể (Trend):} Nhìn bao quát cả giai đoạn 2013--2017, dữ liệu hoàn toàn bị chi phối bởi tính chu kỳ mùa vụ lặp đi lặp lại theo năm. Chưa quan sát thấy một xu hướng giảm bền vững mang tính cấu trúc nào, cho thấy các biện pháp kiểm soát ô nhiễm trong thời kỳ này chưa mang lại hiệu quả rõ rệt đối với khu vực trạm Shunyi.
\end{itemize}

\noindent\textbf{2. Phân phối và biến động theo chu kỳ mùa (Hình \ref{fig:cell7} - bên trái):}
\begin{itemize}
    \item \textbf{Mùa Đông là thời điểm ô nhiễm cực đoan nhất:} Đồ thị Boxplot của mùa Đông có khoảng biến thiên (IQR) rộng nhất và biên trên kéo dài nhất (nhiều thời điểm vượt mốc $300\,\mu\text{g/m}^3$), phản ánh rõ hiệu ứng cộng hưởng từ việc tăng lượng phát thải phục vụ hệ thống sưởi ấm (đốt than) và hiện tượng khí tượng nghịch nhiệt giữ bụi mịn tích tụ sát mặt đất.
    \item \textbf{Mùa Hè có chất lượng không khí tốt nhất:} Nồng độ trung vị (median) và khoảng dao động của mùa Hè đều ở mức thấp nhất nhờ lượng mưa lớn rửa trôi bụi mịn và bức xạ nhiệt mạnh thúc đẩy đối lưu không khí.
    \item \textbf{Mùa Xuân và mùa Thu:} Có mức trung vị khá tương đồng, tuy nhiên mùa Thu bắt đầu xuất hiện dấu hiệu chuyển mùa với biên phân phối phía trên kéo dài hơn, báo hiệu sự khởi đầu của đợt ô nhiễm mùa đông.
\end{itemize}

\noindent\textbf{3. Biến động theo giờ trong ngày và lý giải nghịch lý chất lượng không khí (Hình \ref{fig:cell7} - bên phải):}
Mặc dù giờ cao điểm giao thông và sản xuất diễn ra ban ngày, nồng độ \texttt{PM2.5} sát mặt đất lại ghi nhận mức thấp nhất vào buổi chiều và đạt đỉnh vào đêm muộn. Nghịch lý này được giải thích một cách khoa học qua cơ chế vật lý và đô thị:
\begin{itemize}
    \item \textbf{Biến động lớp biên khí quyển (PBL) và nghịch nhiệt ban đêm:} 
    \begin{itemize}
        \item \textit{Ban ngày (đặc biệt từ 14h--16h):} Bức xạ mặt trời làm nóng mặt đất cực đại, kích hoạt các luồng đối lưu nhiệt mạnh mẽ đẩy chiều cao lớp biên khí quyển lên tới $1.5 - 2.5\,\text{km}$. Thể tích không khí khuếch tán tăng lên gấp nhiều lần, đóng vai trò như một cơ chế pha loãng tự nhiên giúp nồng độ bụi sát mặt đất giảm xuống mức cực tiểu.
        \item \textit{Ban đêm (đặc biệt từ 23h--02h):} Quá trình tỏa nhiệt hồng ngoại của mặt đất làm giảm nhanh nhiệt độ sát đất, hình thành hiện tượng nghịch nhiệt bức xạ và hạ thấp chiều cao lớp biên sụt xuống chỉ còn $100 - 300\,\text{m}$ \cite{gao2015}. Toàn bộ lượng phát thải tích tụ từ chiều tối bị giam giữ trong khoảng không gian chật hẹp sát đất, không thể phát tán thẳng đứng hay đối lưu.
    \end{itemize}
    \item \textbf{Quy định phân luồng xe tải hạng nặng của đô thị Bắc Kinh:} Để kiểm soát ùn tắc và ô nhiễm nội đô ban ngày, chính quyền Bắc Kinh cấm tuyệt đối các phương tiện vận tải lớn chạy bằng động cơ diesel (nguồn phát thải bụi mịn sơ cấp và khí tiền chất lớn nhất) di chuyển trong thành phố từ $06$h$00$ đến $23$h$00$. Việc các đoàn xe tải nặng đồng loạt đi vào thành phố sau $23$h$00$ đêm gặp ngay điều kiện lớp biên khí quyển sụt thấp và ổn định là tác nhân trực tiếp tạo nên đỉnh ô nhiễm PM2.5 cực đại ban đêm \cite{biggart2020}.
    \item \textbf{Độ trễ hóa học trong quá trình tạo bụi thứ cấp (Secondary PM2.5):} Bụi mịn PM2.5 trong đô thị chủ yếu được cấu thành từ bụi thứ cấp (chiếm $60\% - 80\%$) qua các phản ứng hóa học của các khí tiền chất ($\text{SO}_2, \text{NO}_x, \text{NH}_3, \text{VOCs}$). Khói bụi thải ra vào giờ cao điểm sáng ($07$h--$09$h) cần khoảng thời gian trễ từ vài giờ đến nửa ngày để phản ứng quang hóa và chuyển hóa pha khí-hạt (gas-to-particle conversion), khiến đỉnh nồng độ PM2.5 thực tế bị dịch chuyển lùi về đêm muộn \cite{seinfeld2016}.
\end{itemize}

\subsection{Kiểm tra \& Phát hiện các vấn đề dữ liệu}

\subsubsection{Dữ liệu trùng lặp}
Để đảm bảo tính toàn vẹn và tính nhất quán của chuỗi thời gian, nghiên cứu tiến hành kiểm tra dữ liệu trùng lặp theo hai cấp độ:
\begin{itemize}
    \item \textbf{Trùng lặp hoàn toàn:} Kiểm tra các bản ghi giống nhau hoàn toàn trên tất cả các trường dữ liệu. Hiện tượng này thường xảy ra do lỗi ghi nhận lặp của cảm biến hoặc lỗi trong quá trình xuất/nhập cơ sở dữ liệu.
    \item \textbf{Trùng lặp mốc thời gian:} Sử dụng tổ hợp các cột định danh thời gian gồm \texttt{year}, \texttt{month}, \texttt{day}, \texttt{hour} làm khóa duy nhất. Chuỗi thời gian yêu cầu mỗi mốc giờ chỉ tồn tại duy nhất một bản ghi đo lường. Nếu phát hiện trùng lặp mốc thời gian, tính liên tục của chuỗi dữ liệu sẽ bị phá vỡ, gây sai lệch nghiêm trọng cho việc tạo các biến trễ (lag features) và phân tích chuỗi thời gian sau này.
\end{itemize}

Kết quả thực thi từ chương trình được hiển thị ở Hình \ref{fig:cell8}:

\begin{figure}[H]
    \centering
    \includegraphics[width=0.85\textwidth]{cell8.png}
    \caption{Kết quả kiểm tra dữ liệu trùng lặp từ mã nguồn phân tích}
    \label{fig:cell8}
\end{figure}

\noindent\textbf{Đánh giá kết quả:}
Hình \ref{fig:cell8} cho thấy số lượng dòng trùng lặp hoàn toàn và số dòng trùng lặp mốc thời gian đều bằng $0$. Điều này xác nhận tập dữ liệu của trạm quan trắc Shunyi có độ sạch tuyệt đối về mặt cấu trúc thời gian với đúng $35.064$ dòng quan sát phân bổ liên tục không gián đoạn trong 4 năm ($1.461 \text{ ngày} \times 24 \text{ giờ}$). Kết quả này thiết lập một nền tảng vững chắc để thực hiện các bước trích xuất đặc trưng thời gian trong các bước tiếp theo của đồ án.


\subsubsection{Dữ liệu bị khuyết thiếu}
Tiến hành tính toán tỷ lệ khuyết thiếu trên từng cột và trực quan hóa:
\begin{figure}[H]
    \centering
    \includegraphics[width=0.9\textwidth]{cell9.2.png}
    \caption{Tỷ lệ khuyết thiếu của các cột đặc trưng trong tập dữ liệu gốc}
\end{figure}

Bộ dữ liệu chứa tỷ lệ khuyết thiếu dao động từ 0.13\% (khí tượng) đến 6.21\% (\texttt{CO} và \texttt{O3}). Mức khuyết thiếu của biến mục tiêu \texttt{PM2.5} là 2.65\%. Vì tỷ lệ thiếu đều nằm quanh ngưỡng ý nghĩa thống kê ($\ge 5\%$), việc xử lý missing values là bắt buộc.

\subsubsection{Kiểm tra logic vật lý và phân phối chuẩn}

\noindent\textbf{1. Kiểm tra logic vật lý (Giá trị đặc biệt):}
Kiểm tra điều kiện điểm sương không vượt quá nhiệt độ thực tế (\texttt{DEWP} $\le$ \texttt{TEMP}), cũng như tốc độ gió và lượng mưa âm.
\begin{figure}[H]
    \centering
    \includegraphics[width=0.9\textwidth]{cell10.png}
    \caption{Kết quả kiểm tra giá trị đặc biệt và mâu thuẫn vật lý từ mã nguồn}
\end{figure}
Kết quả kiểm tra xác nhận hoàn toàn không có bất kỳ dòng dữ liệu nào xảy ra mâu thuẫn vật lý (\texttt{DEWP} $>$ \texttt{TEMP}), tức là nhiệt độ điểm sương luôn nhỏ hơn hoặc bằng nhiệt độ không khí thực tế. Điều này minh chứng cho tính hợp lệ vật lý và độ tin cậy khí tượng rất cao của dữ liệu đo đạc tại trạm Shunyi. Mặc dù vậy, để tăng cường tính chống chịu lỗi (fault-tolerance), pipeline tiền xử lý vẫn tích hợp bước kiểm tra logic này (nếu xuất hiện mâu thuẫn sẽ tự động hiệu chỉnh gán \texttt{DEWP} = \texttt{TEMP}) như một cơ chế phòng ngừa rủi ro cho luồng dữ liệu.


\noindent\textbf{2. Kiểm định phân phối chuẩn:}
Sử dụng kiểm định D'Agostino-Pearson \cite{dagostino1973} bác bỏ giả thuyết $H_0$ (phân phối chuẩn) cho tất cả các chất ô nhiễm ($p < 0.001$), xác nhận sự cần thiết của phép biến đổi logarit để ổn định phương sai sai số.
\begin{figure}[H]
    \centering
    \includegraphics[width=0.9\textwidth]{cell11.png}
    \caption{Kết quả kiểm tra phân phối chuẩn (D'Agostino-Pearson test) từ mã nguồn}
\end{figure}
Như hình trên thể hiện, toàn bộ các biến ô nhiễm không đạt phân phối chuẩn, trong khi các biến khí tượng gần chuẩn hơn.

\noindent\textbf{3. Kiểm tra đa cộng tuyến ban đầu:}
Kết quả phân tích VIF từ mã nguồn (Hình \ref{fig:cell12}) chỉ ra rằng mặc dù không có biến độc lập nào vượt quá ngưỡng đa cộng tuyến cực độ ($VIF > 10$), các biến thuộc hai cặp đặc trưng tương đồng gồm \texttt{TEMP} ($8.37$) -- \texttt{DEWP} ($6.14$) và \texttt{PM2.5} ($7.63$) -- \texttt{PM10} ($6.12$) đều vượt quá ngưỡng cảnh báo đa cộng tuyến trung bình ($VIF > 5$). Điều này phản ánh sự tồn tại của hiện tượng đa cộng tuyến đáng kể giữa các đặc trưng nhiệt động (nhiệt độ thực tế và nhiệt độ điểm sương) cũng như giữa các hạt bụi mịn có chung nguồn phát và cơ chế khí quyển.
\begin{figure}[H]
    \centering
    \includegraphics[width=0.9\textwidth]{cell12.png}
    \caption{Kết quả phân tích chỉ số VIF ban đầu từ mã nguồn}
    \label{fig:cell12}
\end{figure}
Hiện tượng đa cộng tuyến trung bình này vẫn có thể làm phình to phương sai của các hệ số ước lượng, làm giảm độ tin cậy của các kiểm định thống kê và gây mất ổn định số học khi giải hệ phương trình chuẩn OLS. Hướng khắc phục tối ưu là loại bỏ đặc trưng phụ thuộc tuyến tính trực tiếp \texttt{PM10} (chỉ giữ lại \texttt{PM2.5} làm biến mục tiêu hồi quy) và áp dụng phân rã QR có chọn trục xoay (Pivoting QR Decomposition) trong pipeline để loại bỏ hoàn toàn các cột phụ thuộc tuyến tính trong ma trận thiết kế ở các bước tiếp theo.


\subsection{Xử lý Dữ liệu khuyết thiếu}

\subsubsection{Phân tích cơ chế thiếu}
\begin{figure}[H]
    \centering
    \includegraphics[width=0.9\textwidth]{cell14.png}
    \caption{Bản đồ trực quan hóa cụm khuyết thiếu theo thời gian}
\end{figure}

Phân tích bản đồ khuyết thiếu cho thấy:
\begin{itemize}
    \item \textbf{Các đặc trưng khí tượng (\texttt{TEMP}, \texttt{PRES}, \texttt{DEWP}, \texttt{RAIN}, \texttt{WSPM}):} Tỷ lệ thiếu cực thấp ($< 0.15\%$), xuất hiện ngẫu nhiên đơn lẻ $\rightarrow$ Cơ chế khuyết thiếu hoàn toàn ngẫu nhiên MCAR.
    \item \textbf{Các chất ô nhiễm (\texttt{PM2.5}, v.v.):} Thiếu tập trung theo từng cụm thời gian dài liên tục (do trạm quan trắc mất điện hoặc bảo trì cảm biến định kỳ) $\rightarrow$ Cơ chế khuyết thiếu ngẫu nhiên MAR.
    \item \textbf{Hướng gió (\texttt{wd}):} Thường khuyết thiếu đồng thời khi tốc độ gió bằng 0 (trời lặng gió, cảm biến hướng gió không xoay được) $\rightarrow$ Cơ chế khuyết thiếu ngẫu nhiên MAR.
\end{itemize}


\subsubsection{Lý do chọn phương pháp xử lý và hiện thực hóa}
Dựa trên cơ chế thiếu toán học ở trên, các phương pháp tối ưu được lựa chọn cụ thể:
\begin{itemize}
    \item \textbf{Khí tượng MCAR $\rightarrow$ Nội suy tuyến tính theo thời gian:} Vì tỷ lệ khuyết thiếu cực kỳ nhỏ và đặc trưng thời tiết biến thiên liên tục theo thời gian, nội suy tuyến tính là phương án tự nhiên và mượt mà nhất.
    \item \textbf{Chất ô nhiễm MAR $\rightarrow$ Phương pháp KNN Imputation (k=5):} Việc gán bằng trung bình hoặc trung vị sẽ làm triệt tiêu phương sai và bóp méo phân phối thực tế. KNN Imputation tận dụng mối tương quan phi tuyến đa biến mạnh mẽ giữa các chất ô nhiễm và nhiệt độ tại cùng thời điểm để khôi phục giá trị khuyết thiếu một cách chân thực nhất.
    \item \textbf{Hướng gió \texttt{wd} $\rightarrow$ Mode theo cửa sổ di động 24 giờ:} Là biến phân loại định tính. Hướng gió có xu hướng giữ nguyên hoặc lặp lại theo chu kỳ ngày/đêm ngắn, do đó ta tính giá trị xuất hiện nhiều nhất (Mode) trong vòng $\pm 12$ giờ xung quanh điểm khuyết thiếu để điền dữ liệu.
\end{itemize}

\subsection{Pipeline Tiền Xử Lý Tổng Thể}

\subsubsection{Sơ đồ luồng xử lý toán học của Pipeline}
Để đảm bảo mô hình hồi quy tuyến tính hoạt động hiệu quả nhất và tuyệt đối không xảy ra hiện tượng rò rỉ dữ liệu từ tập Test sang tập Train, ta đóng gói toàn bộ luồng xử lý vào một class \textbf{DataPipeline} chuyên nghiệp với cấu trúc \textbf{fit} và \textbf{transform} chuẩn mực:

\begin{center}
\begin{tcolorbox}[
    enhanced,
    colback=HCMUSBlue,
    colframe=HCMUSBlue,
    coltitle=white,
    fonttitle=\bfseries\large,
    halign=center,
    valign=center,
    arc=4pt,
    boxrule=1pt,
    width=0.7\textwidth
]
    \centering
    \bfseries \color{white} DỮ LIỆU THÔ (RAW DATA)\\
    \textnormal{\small 35.064 dòng $\times$ 18 cột}
\end{tcolorbox}
\end{center}

\begin{center}
    \tikz[baseline=-0.5ex]\draw[->,line width=2pt,HCMUSBlue] (0,0.3) -- (0,-0.3);
\end{center}

\begin{tcolorbox}[
    enhanced,
    colback=white,
    colframe=HCMUSBlue,
    colbacktitle=HCMUSBlue!8,
    coltitle=HCMUSBlue,
    fonttitle=\bfseries,
    title={Bước 1: Làm sạch ban đầu},
    attach boxed title to top left={yshift=-2mm,xshift=3mm},
    boxed title style={colback=HCMUSBlue!8,arc=3pt,boxrule=0.8pt,colframe=HCMUSBlue},
    boxrule=0.8pt,
    arc=4pt,
    left=12pt,
    right=12pt,
    top=10pt,
    bottom=8pt,
    boxsep=2pt,
    before=\vspace{1mm},
    after=\vspace{1mm}
]
    \begin{itemize}[leftmargin=*,itemsep=3pt,topsep=0pt,parsep=0pt]
        \item \textbf{Loại bỏ đặc trưng không liên quan:} Loại bỏ biến rác \texttt{'No'} (số thứ tự) và \texttt{'station'} (tên trạm - do chỉ có 1 trạm duy nhất).
        \item \textbf{Hiệu chỉnh logic vật lý:} Gán \texttt{DEWP} = \texttt{TEMP} nếu xảy ra mâu thuẫn \texttt{DEWP} $>$ \texttt{TEMP} (cơ chế phòng ngừa rủi ro).
        \item \textbf{Giải quyết đa cộng tuyến từ gốc:} Loại bỏ đặc trưng phụ thuộc tuyến tính trực tiếp \texttt{'PM10'} để chỉ tập trung dự báo biến mục tiêu \texttt{'PM2.5'}.
        \item \textbf{Giới hạn dưới vật lý:} Cắt các giá trị âm của các chất ô nhiễm về 0 (\texttt{clip(lower=0)}) để tránh nồng độ âm phi thực tế.
    \end{itemize}
\end{tcolorbox}

\begin{center}
    \tikz[baseline=-0.5ex]\draw[->,line width=2pt,HCMUSBlue] (0,0.3) -- (0,-0.3);
\end{center}

\begin{tcolorbox}[
    enhanced,
    colback=white,
    colframe=HCMUSBlue,
    colbacktitle=HCMUSBlue!8,
    coltitle=HCMUSBlue,
    fonttitle=\bfseries,
    title={Bước 2: Điền khuyết thiếu (Imputation)},
    attach boxed title to top left={yshift=-2mm,xshift=3mm},
    boxed title style={colback=HCMUSBlue!8,arc=3pt,boxrule=0.8pt,colframe=HCMUSBlue},
    boxrule=0.8pt,
    arc=4pt,
    left=12pt,
    right=12pt,
    top=10pt,
    bottom=8pt,
    boxsep=2pt,
    before=\vspace{1mm},
    after=\vspace{1mm}
]
    \begin{itemize}[leftmargin=*,itemsep=3pt,topsep=0pt,parsep=0pt]
        \item \textbf{Nội suy tuyến tính (Linear Interpolation):} Áp dụng cho các biến khí tượng định lượng (tỷ lệ thiếu cực thấp, biến thiên liên tục theo chuỗi thời gian).
        \item \textbf{Mode theo cửa sổ di động 24 giờ:} Áp dụng điền hướng gió định tính \texttt{'wd'} dựa trên tần suất xuất hiện nhiều nhất trong cửa sổ $\pm 12$ giờ.
        \item \textbf{KNN Imputer (k=5):} Điền khuyết thiếu cho biến mục tiêu \texttt{'PM2.5'} và các chất ô nhiễm dựa trên mối tương quan phi tuyến đa biến mạnh mẽ với các trạm lân cận và đặc trưng khí tượng.
    \end{itemize}
\end{tcolorbox}

\begin{center}
    \tikz[baseline=-0.5ex]\draw[->,line width=2pt,HCMUSBlue] (0,0.3) -- (0,-0.3);
\end{center}

\begin{tcolorbox}[
    enhanced,
    colback=white,
    colframe=HCMUSBlue,
    colbacktitle=HCMUSBlue!8,
    coltitle=HCMUSBlue,
    fonttitle=\bfseries,
    title={Bước 3: Tách dữ liệu Train/Test theo thời gian},
    attach boxed title to top left={yshift=-2mm,xshift=3mm},
    boxed title style={colback=HCMUSBlue!8,arc=3pt,boxrule=0.8pt,colframe=HCMUSBlue},
    boxrule=0.8pt,
    arc=4pt,
    left=12pt,
    right=12pt,
    top=10pt,
    bottom=8pt,
    boxsep=2pt,
    before=\vspace{1mm},
    after=\vspace{1mm}
]
    \begin{itemize}[leftmargin=*,itemsep=3pt,topsep=0pt,parsep=0pt]
        \item \textbf{Nguyên tắc phân rã thời gian:} Tránh rò rỉ thông tin tương lai sang quá khứ.
        \item \textbf{Tập huấn luyện (Train Set):} Trước ngày 01/09/2016 (bao gồm 30,720 mẫu, chiếm khoảng $87.6\%$).
        \item \textbf{Tập kiểm thử (Test Set):} Từ ngày 01/09/2016 trở đi (bao gồm 4,344 mẫu, chiếm khoảng $12.4\%$).
    \end{itemize}
\end{tcolorbox}

\begin{center}
    \tikz[baseline=-0.5ex]\draw[->,line width=2pt,HCMUSBlue] (0,0.3) -- (0,-0.3);
\end{center}

\begin{tcolorbox}[
    enhanced,
    colback=white,
    colframe=HCMUSBlue,
    colbacktitle=HCMUSBlue!8,
    coltitle=HCMUSBlue,
    fonttitle=\bfseries,
    title={Bước 4: Loại bỏ phần đuôi nặng (Outliers)},
    attach boxed title to top left={yshift=-2mm,xshift=3mm},
    boxed title style={colback=HCMUSBlue!8,arc=3pt,boxrule=0.8pt,colframe=HCMUSBlue},
    boxrule=0.8pt,
    arc=4pt,
    left=12pt,
    right=12pt,
    top=10pt,
    bottom=8pt,
    boxsep=2pt,
    before=\vspace{1mm},
    after=\vspace{1mm}
]
    \begin{itemize}[leftmargin=*,itemsep=3pt,topsep=0pt,parsep=0pt]
        \item \textbf{Phương pháp Winsorization:} Áp dụng giới hạn trên đối với các biến chất ô nhiễm ở tập huấn luyện (Train Set).
        \item \textbf{Ngưỡng cắt lọc:} Giới hạn giá trị lớn nhất tại mức $Q_3 + 1.5 \times \text{IQR}$ của tập Train để triệt tiêu ảnh hưởng của các điểm đo dị biệt cực đoan làm méo mó mô hình hồi quy.
    \end{itemize}
\end{tcolorbox}

\begin{center}
    \tikz[baseline=-0.5ex]\draw[->,line width=2pt,HCMUSBlue] (0,0.3) -- (0,-0.3);
\end{center}

\begin{tcolorbox}[
    enhanced,
    colback=white,
    colframe=HCMUSBlue,
    colbacktitle=HCMUSBlue!8,
    coltitle=HCMUSBlue,
    fonttitle=\bfseries,
    title={Bước 5: Biến đổi ổn định phương sai},
    attach boxed title to top left={yshift=-2mm,xshift=3mm},
    boxed title style={colback=HCMUSBlue!8,arc=3pt,boxrule=0.8pt,colframe=HCMUSBlue},
    boxrule=0.8pt,
    arc=4pt,
    left=12pt,
    right=12pt,
    top=10pt,
    bottom=8pt,
    boxsep=2pt,
    before=\vspace{1mm},
    after=\vspace{1mm}
]
    \begin{itemize}[leftmargin=*,itemsep=3pt,topsep=0pt,parsep=0pt]
        \item \textbf{Phép biến đổi Logarithm tự nhiên (Log1p):} Áp dụng $\text{log\_PM2.5} = \ln(\text{PM2.5} + 1)$ cho biến mục tiêu.
        \item \textbf{Ý nghĩa thống kê:} Đưa phân phối lệch phải nghiêm trọng về gần phân phối chuẩn, ổn định phương sai sai số và tuyến tính hóa mối quan hệ với các đặc trưng khí tượng.
    \end{itemize}
\end{tcolorbox}

\begin{center}
    \tikz[baseline=-0.5ex]\draw[->,line width=2pt,HCMUSBlue] (0,0.3) -- (0,-0.3);
\end{center}

\begin{tcolorbox}[
    enhanced,
    colback=white,
    colframe=HCMUSBlue,
    colbacktitle=HCMUSBlue!8,
    coltitle=HCMUSBlue,
    fonttitle=\bfseries,
    title={Bước 6: Tái cấu trúc đặc trưng thời gian (Gom nhóm - Binning)},
    attach boxed title to top left={yshift=-2mm,xshift=3mm},
    boxed title style={colback=HCMUSBlue!8,arc=3pt,boxrule=0.8pt,colframe=HCMUSBlue},
    boxrule=0.8pt,
    arc=4pt,
    left=12pt,
    right=12pt,
    top=10pt,
    bottom=8pt,
    boxsep=2pt,
    before=\vspace{1mm},
    after=\vspace{1mm}
]
    \begin{itemize}[leftmargin=*,itemsep=3pt,topsep=0pt,parsep=0pt]
        \item \textbf{Biến \texttt{'month'} $\rightarrow$ 4 Mùa (season):} Xuân, Hạ, Thu, Đông (nắm bắt tính chất mùa vụ sâu sắc của khí hậu Bắc Kinh).
        \item \textbf{Biến \texttt{'hour'} $\rightarrow$ 4 Buổi (time\_of\_day):} Sáng, Trưa, Chiều, Tối (phục vụ lý giải cơ chế lớp biên khí quyển).
        \item \textbf{Biến \texttt{'day'} $\rightarrow$ 7 Thứ trong tuần (day\_of\_week):} Đại diện cho chu kỳ sinh hoạt đô thị và mật độ giao thông ngày thường so với ngày nghỉ.
    \end{itemize}
\end{tcolorbox}

\begin{center}
    \tikz[baseline=-0.5ex]\draw[->,line width=2pt,HCMUSBlue] (0,0.3) -- (0,-0.3);
\end{center}

\begin{tcolorbox}[
    enhanced,
    colback=white,
    colframe=HCMUSBlue,
    colbacktitle=HCMUSBlue!8,
    coltitle=HCMUSBlue,
    fonttitle=\bfseries,
    title={Bước 7: Chế tạo đặc trưng nâng cao (Feature Engineering)},
    attach boxed title to top left={yshift=-2mm,xshift=3mm},
    boxed title style={colback=HCMUSBlue!8,arc=3pt,boxrule=0.8pt,colframe=HCMUSBlue},
    boxrule=0.8pt,
    arc=4pt,
    left=12pt,
    right=12pt,
    top=10pt,
    bottom=8pt,
    boxsep=2pt,
    before=\vspace{1mm},
    after=\vspace{1mm}
]
    \begin{itemize}[leftmargin=*,itemsep=3pt,topsep=0pt,parsep=0pt]
        \item \textbf{Biến trễ tự hồi quy (Time-Lags):} Tạo \texttt{log\_PM2.5\_lag\_1}, \texttt{log\_PM2.5\_lag\_2}, \texttt{log\_PM2.5\_lag\_24} để thu thập quán tính thời gian mạnh mẽ của chất lượng không khí.
        \item \textbf{Biến tương tác phi tuyến (Interactions):} \texttt{TEMP} $\times$ \texttt{WSPM} (nhiệt độ kết hợp tốc độ gió) và \texttt{RAIN} $\times$ \texttt{WSPM} (lượng mưa kết hợp tốc độ gió).
        \item \textbf{Biến đa thức (Polynomials):} \texttt{TEMP\_sq} (nhiệt độ bình phương), \texttt{WSPM\_sq} (tốc độ gió bình phương) để mô hình hóa các biến động phi tuyến vật lý.
    \end{itemize}
\end{tcolorbox}

\begin{center}
    \tikz[baseline=-0.5ex]\draw[->,line width=2pt,HCMUSBlue] (0,0.3) -- (0,-0.3);
\end{center}

\begin{tcolorbox}[
    enhanced,
    colback=white,
    colframe=HCMUSBlue,
    colbacktitle=HCMUSBlue!8,
    coltitle=HCMUSBlue,
    fonttitle=\bfseries,
    title={Bước 8: Mã hóa \& Chuẩn hóa (OHE \& Scaling)},
    attach boxed title to top left={yshift=-2mm,xshift=3mm},
    boxed title style={colback=HCMUSBlue!8,arc=3pt,boxrule=0.8pt,colframe=HCMUSBlue},
    boxrule=0.8pt,
    arc=4pt,
    left=12pt,
    right=12pt,
    top=10pt,
    bottom=8pt,
    boxsep=2pt,
    before=\vspace{1mm},
    after=\vspace{1mm}
]
    \begin{itemize}[leftmargin=*,itemsep=3pt,topsep=0pt,parsep=0pt]
        \item \textbf{One-Hot Encoding (OHE):} Áp dụng cho toàn bộ các biến phân loại định tính vừa nhắc đến ở Bước 6.
        \item \textbf{Chuẩn hóa StandardScaler:} Thực hiện chuẩn hóa chuẩn tắc ($z$-score scaling) với nguyên tắc nghiêm ngặt: chỉ học trung bình $\mu$ và độ lệch chuẩn $\sigma$ trên tập huấn luyện (Train Set), sau đó áp dụng biến đổi (\texttt{transform}) lên cả tập Train và tập Test.
    \end{itemize}
\end{tcolorbox}

\begin{center}
    \tikz[baseline=-0.5ex]\draw[->,line width=2pt,HCMUSBlue] (0,0.3) -- (0,-0.3);
\end{center}

\begin{tcolorbox}[
    enhanced,
    colback=white,
    colframe=HCMUSBlue,
    colbacktitle=HCMUSBlue!8,
    coltitle=HCMUSBlue,
    fonttitle=\bfseries,
    title={Bước 9: Phân rã ma trận QR (QR Decomposition có Pivoting)},
    attach boxed title to top left={yshift=-2mm,xshift=3mm},
    boxed title style={colback=HCMUSBlue!8,arc=3pt,boxrule=0.8pt,colframe=HCMUSBlue},
    boxrule=0.8pt,
    arc=4pt,
    left=12pt,
    right=12pt,
    top=10pt,
    bottom=8pt,
    boxsep=2pt,
    before=\vspace{1mm},
    after=\vspace{1mm}
]
    \begin{itemize}[leftmargin=*,itemsep=3pt,topsep=0pt,parsep=0pt]
        \item \textbf{Phát hiện hạng ma trận tự động:} Thực hiện phân rã QR có chọn trục xoay (Pivoting QR) để tự động phát hiện hạng cột (column rank) của ma trận đặc trưng mở rộng $X$.
        \item \textbf{Loại bỏ cột phụ thuộc tuyến tính hoàn hảo:} Loại bỏ triệt để các cột dư thừa/phụ thuộc tuyến tính hoàn hảo phát sinh do mã hóa OHE và tương tác, chống hiện tượng suy biến ma trận (Singular Matrix) giúp phép toán bình phương tối thiểu luôn ổn định và hội tụ.
    \end{itemize}
\end{tcolorbox}

\begin{center}
    \tikz[baseline=-0.5ex]\draw[->,line width=2pt,HCMUSBlue] (0,0.3) -- (0,-0.3);
\end{center}

\begin{center}
\begin{tcolorbox}[
    enhanced,
    colback=CommentGreen!10,
    colframe=CommentGreen,
    coltitle=CommentGreen,
    fonttitle=\bfseries\large,
    halign=center,
    valign=center,
    arc=4pt,
    boxrule=1pt,
    width=0.7\textwidth
]
    \centering
    \bfseries \color{CommentGreen!60!black} SẴN SÀNG HUẤN LUYỆN HỒI QUY\\
    \textnormal{\small Ma trận thiết kế sạch tuyệt đối với đúng 40 đặc trưng chất lượng}
\end{tcolorbox}
\end{center}


\subsubsection{Phân rã QR loại bỏ đa cộng tuyến hoàn hảo}
Quá trình tạo biến đa thức và One-hot encoding vô tình tạo ra sự phụ thuộc tuyến tính tuyệt đối giữa các cột đặc trưng mới (định thức ma trận $X^T X$ bằng 0 hoàn toàn), làm suy biến ma trận thiết kế và gây ra lỗi toán học nghiêm trọng (Singular Matrix) khi giải hệ phương trình chuẩn OLS.

Để khắc phục triệt để, ta áp dụng \textbf{Phân rã QR có chọn phần tử chốt} lên ma trận thiết kế $X$ trước khi đưa vào huấn luyện:
\begin{equation}
    X P = Q R
\end{equation}
Trong đó:
\begin{itemize}
    \item $P$ là ma trận hoán vị cột.
    \item $Q \in \mathbb{R}^{n \times n}$ là ma trận trực giao ($Q^T Q = I$).
    \item $R \in \mathbb{R}^{n \times (p+1)}$ là ma trận tam giác trên có các phần tử đường chéo $|r_{11}| \ge |r_{22}| \ge \dots \ge |r_{kk}| \ge 0$.
\end{itemize} Kỹ thuật này 
Bằng cách theo dõi các phần tử đường chéo của $R$, tại vị trí phần tử đường chéo $|r_{jj}| < 10^{-10}$ (xấp xỉ bằng 0), thuật toán xác định cột thứ $j$ là cột phụ thuộc tuyến tính tuyệt đối vào các cột trước đó và thẳng tay loại bỏ cột này. Kỹ thuật này đảm bảo ma trận thiết kế cuối cùng luôn đạt hạng đầy đủ cột một cách hoàn hảo về mặt toán học.

\subsection{Xây dựng các Mô hình Hồi quy và Giải thuật tự cài đặt}
Để giải quyết bài toán dự báo chất lượng không khí, đồ án tiến hành thiết kế và cài đặt hoàn toàn từ đầu (from scratch) bằng ngôn ngữ Python sử dụng thư viện NumPy 6 mô hình hồi quy khác nhau. Các thuật toán tối ưu hóa, từ nghiệm đóng của phương trình chuẩn tắc đến các thuật toán lặp và hồi quy phi tuyến, đều được tự hiện thực hóa để đảm bảo tính minh bạch toán học và độ ổn định cao nhất.

\subsubsection{Mô hình 1: OLS cơ bản (Ordinary Least Squares)}
OLS tìm kiếm vector tham số $\boldsymbol{\beta}$ sao cho tối thiểu hóa tổng bình phương sai số (RSS):
\begin{equation}
    \hat{\boldsymbol{\beta}}_{\text{OLS}} = \arg\min_{\boldsymbol{\beta}} \| \mathbf{y} - \mathbf{X}\boldsymbol{\beta} \|^2
\end{equation}
Nghiệm đóng của bài toán được xác định qua phương trình chuẩn tắc (Normal Equation):
\begin{equation}
    \hat{\boldsymbol{\beta}}_{\text{OLS}} = (\mathbf{X}^\top\mathbf{X})^{-1} \mathbf{X}^\top\mathbf{y}
\end{equation}
\textbf{Hiện thực hóa thuật toán:} Để tránh các vấn đề không ổn định số học khi nghịch đảo trực tiếp ma trận $(\mathbf{X}^\top\mathbf{X})^{-1}$, chương trình sử dụng phương pháp giải hệ phương trình bằng thuật toán phân rã LU thông qua hàm \texttt{np.linalg.solve}, giúp đạt tốc độ và độ chính xác tối ưu.

\subsubsection{Mô hình 2: OLS với chọn biến ngược (Backward Elimination)}
Mô hình OLS Basic sử dụng toàn bộ 40 đặc trưng. Tuy nhiên, một số đặc trưng có thể không có ý nghĩa thống kê hoặc gây ra đa cộng tuyến nghiêm trọng. Kỹ thuật loại bỏ ngược (Backward Elimination) được tự động hóa dựa trên hai chỉ số:
\begin{enumerate}
    \item \textbf{Variance Inflation Factor (VIF):} Đo lường mức độ đa cộng tuyến của biến thứ $j$:
    \begin{equation}
        \text{VIF}(j) = \frac{1}{1 - R^2_j}
    \end{equation}
    với $R^2_j$ là hệ số xác định khi hồi quy đặc trưng $X_j$ theo toàn bộ các đặc trưng còn lại.
    \item \textbf{Trị số p-value của kiểm định t:} Đánh giá ý nghĩa thống kê của từng hệ số hồi quy $\beta_j$.
\end{enumerate}
\textbf{Giải thuật loại bỏ ngược:}
\begin{itemize}
    \item \textit{Bước 1:} Huấn luyện mô hình OLS trên tập biến hiện tại.
    \item \textit{Bước 2:} Tính toán p-value và VIF cho từng đặc trưng.
    \item \textit{Bước 3:} Ưu tiên kiểm tra đa cộng tuyến: Nếu đặc trưng có VIF lớn nhất vượt ngưỡng $\text{VIF}_{\text{threshold}} = 10.0$, loại bỏ đặc trưng này.
    \item \textit{Bước 4:} Nếu không có đa cộng tuyến nghiêm trọng, tìm đặc trưng có p-value lớn nhất. Nếu p-value lớn hơn mức ý nghĩa $\alpha = 0.05$, tiến hành loại bỏ đặc trưng đó.
    \item \textit{Bước 5:} Lặp lại quy trình cho đến khi tất cả các đặc trưng còn lại đều vượt qua cả hai ngưỡng kiểm tra.
\end{itemize}

\subsubsection{Mô hình 3: Hồi quy Ridge (Regularization L2)}
Để hạn chế hiện tượng quá khớp (overfitting) và ổn định các ước lượng hệ số khi ma trận thiết kế bị đa cộng tuyến, mô hình Ridge thêm một thành phần phạt L2 (bình phương chuẩn Euclid) vào hàm mất mát:
\begin{equation}
    \hat{\boldsymbol{\beta}}_{\text{Ridge}} = \arg\min_{\boldsymbol{\beta}} \| \mathbf{y} - \mathbf{X}\boldsymbol{\beta} \|^2 + \lambda \|\boldsymbol{\beta}\|_2^2
\end{equation}
Nghiệm đóng tối ưu của Ridge Regression được tính bằng:
\begin{equation}
    \hat{\boldsymbol{\beta}}_{\text{Ridge}} = (\mathbf{X}^\top\mathbf{X} + \lambda \mathbf{I}_{\text{mod}})^{-1} \mathbf{X}^\top\mathbf{y}
\end{equation}
Trong đó $\mathbf{I}_{\text{mod}}$ là ma trận đơn vị hiệu chỉnh với phần tử đường chéo đầu tiên $I_{00} = 0$, đảm bảo \textbf{không phạt hệ số chặn (intercept)} do hệ số chặn chỉ đóng vai trò dịch chuyển vị trí siêu phẳng mà không gây ra overfitting.

\subsubsection{Mô hình 4: Hồi quy Lasso (Regularization L1)}
Khác với Ridge, Lasso sử dụng thành phần phạt L1 (trị tuyệt đối) giúp thu hẹp các hệ số không quan trọng về \textbf{đúng bằng 0}, đóng vai trò lựa chọn đặc trưng tự động:
\begin{equation}
    \hat{\boldsymbol{\beta}}_{\text{Lasso}} = \arg\min_{\boldsymbol{\beta}} \frac{1}{2n}\| \mathbf{y} - \mathbf{X}\boldsymbol{\beta} \|^2 + \lambda \|\boldsymbol{\beta}\|_1
\end{equation}
Do hàm mất mát của Lasso không khả vi tại điểm $\beta_j = 0$, không tồn tại nghiệm đóng. Thuật toán \textbf{Coordinate Descent} (Tối ưu hóa từng tọa độ) được tự cài đặt để giải quyết bài toán lặp này.
\textbf{Thuật toán Coordinate Descent từ đầu:}
Với mỗi chu kỳ lặp, ta tối ưu lần lượt từng hệ số $\beta_j$ trong khi giữ nguyên các hệ số còn lại:
\begin{enumerate}
    \item Tính phần dư bộ phận (partial residual):
    \begin{equation}
        \mathbf{r}_j = \mathbf{y} - \mathbf{X}\boldsymbol{\beta} + \mathbf{X}_j \beta_j
    \end{equation}
    \item Tính hệ số tương quan chưa phạt:
    \begin{equation}
        z_j = \frac{1}{n} \mathbf{X}_j^\top \mathbf{r}_j
    \end{equation}
    \item Áp dụng toán tử ngưỡng mềm (Soft-Thresholding Operator) để xác định giá trị mới của $\beta_j$:
    \begin{equation}
        \beta_j \leftarrow \frac{S(z_j, \lambda)}{\|\mathbf{X}_j\|^2 / n}
    \end{equation}
    với hàm ngưỡng mềm $S(z, \gamma) = \text{sign}(z) \cdot \max(|z| - \gamma, 0)$.
\end{enumerate}
Vòng lặp tiếp tục cho đến khi mức độ thay đổi tối đa giữa các hệ số $\max_j |\beta_j^{\text{new}} - \beta_j^{\text{old}}| < 10^{-4}$ hoặc đạt số vòng lặp tối đa 2000. Hệ số chặn $j=0$ được miễn phạt L1 hoàn toàn.

\subsubsection{Mô hình 5: Hồi quy Kernel Ridge (Advanced Non-Linear)}
Hồi quy Kernel Ridge (KRR) kết hợp mô hình Ridge với kỹ thuật hạt nhân (Kernel Trick) để ánh xạ dữ liệu đầu vào lên không gian đặc trưng vô hạn chiều, cho phép biểu diễn các quan hệ phi tuyến phức tạp \cite{bishop}:
\begin{equation}
    \hat{\mathbf{y}}(x) = \mathbf{k}(x)^\top (\mathbf{K} + \lambda \mathbf{I})^{-1} \mathbf{y}
\end{equation}
Trong đó $\mathbf{K} \in \mathbb{R}^{n_s \times n_s}$ là ma trận Gram hạt nhân với các phần tử được tính toán bằng hàm nhân RBF (Gaussian Kernel):
\begin{equation}
    K_{ij} = \exp\left(-\frac{\|\mathbf{x}_i - \mathbf{x}_j\|^2}{2\ell^2}\right)
\end{equation}
với độ dài đặc trưng $\ell = 10.0$ và hệ số điều chuẩn $\lambda = 0.1$.
\textbf{Tối ưu hóa bộ nhớ thực tế:} Vì tập dữ liệu huấn luyện thực tế rất lớn ($n > 30,000$), việc tính toán và nghịch đảo ma trận Gram có kích thước $n \times n$ sẽ tiêu tốn bộ nhớ RAM khổng lồ ($\approx 7.2$ GB) và có thể gây tràn bộ nhớ. Để khắc phục, thuật toán tự động áp dụng cơ chế subsampling giới hạn số lượng mẫu huấn luyện tối đa là $n_s = 3000$ mẫu phân bố ngẫu nhiên đại diện, đảm bảo mô hình chạy ổn định, nhanh chóng mà vẫn nắm bắt tốt cấu trúc phi tuyến của phân phối.

\subsubsection{Mô hình 6: Hồi quy Bayesian Linear (Probabilistic Modeling)}
Bayesian Linear Regression tiếp cận bài toán hồi quy dưới góc nhìn xác suất, cung cấp không chỉ dự báo điểm mà còn định lượng mức độ không chắc chắn (Uncertainty Quantification) của dự báo thông qua phân phối tiên nghiệm và hậu nghiệm \cite{murphy}.
Ta giả định phân phối tiên nghiệm chuẩn (conjugate normal prior) của tham số $\boldsymbol{\beta} \sim \mathcal{N}(0, \alpha^{-1} \mathbf{I})$ với độ chính xác tiên nghiệm $\alpha = 1.0$.
\textbf{Giải thuật ước lượng phương sai nhiễu hai bước (Two-Pass Noise Estimation):}
Do phương sai của sai số nhiễu $\sigma^2$ chưa biết, thuật toán được cài đặt theo quy trình hai bước để ước lượng $\sigma^2$ chính xác từ thực nghiệm:
\begin{itemize}
    \item \textit{Bước 1:} Khởi tạo mô hình với ước lượng nhiễu ban đầu $\sigma^2 = 1.0$. Tính toán phân phối hậu nghiệm tạm thời và tìm giá trị trung bình hậu nghiệm $m_n$.
    \item \textit{Bước 2:} Tính toán phần dư thực tế trên tập huấn luyện dựa trên $m_n$, từ đó ước lượng không chệch phương sai nhiễu thực tế:
    \begin{equation}
        \hat{\sigma}^2 = \frac{\sum_{i=1}^n (y_i - \mathbf{x}_{b, i}^\top m_n)^2}{n - p}
    \end{equation}
    \item \textit{Bước 3:} Cập nhật lại phân phối hậu nghiệm chính thức bằng phương sai nhiễu $\hat{\sigma}^2$ vừa tìm được:
    \begin{align}
        \mathbf{S}_n &= \left( \alpha \mathbf{I} + \frac{1}{\hat{\sigma}^2} \mathbf{X}_b^\top \mathbf{X}_b \right)^{-1} \\
        \mathbf{m}_n &= \mathbf{S}_n \left( \frac{1}{\hat{\sigma}^2} \mathbf{X}_b^\top \mathbf{y} \right)
    \end{align}
\end{itemize}
Dự báo MAP (Maximum A Posteriori) cho điểm dữ liệu mới là $\hat{y} = \mathbf{x}_b^\top \mathbf{m}_n$. 
Độ lệch chuẩn dự báo hậu nghiệm (posterior predictive standard deviation) dùng để định lượng sai số không chắc chắn được xác định bởi:
\begin{equation}
    \sigma_{\text{pred}}(\mathbf{x}) = \sqrt{\hat{\sigma}^2 + \mathbf{x}_b^\top \mathbf{S}_n \mathbf{x}_b}
\end{equation}

\subsubsection{Quy trình chọn siêu tham số bằng 5-Fold Cross-Validation}
Với các mô hình hồi quy điều chuẩn Ridge và Lasso, giá trị siêu tham số $\lambda$ có vai trò quyết định hiệu năng của mô hình. Để lựa chọn $\lambda$ tối ưu một cách khách quan nhất, quy trình K-Fold Cross-Validation với $K = 5$ được tự cài đặt theo nguyên tắc:
\begin{enumerate}
    \item Chia tập dữ liệu huấn luyện ngẫu nhiên thành 5 phần (folds) bằng nhau.
    \item Thiết lập danh sách 50 ứng viên siêu tham số $\lambda$ phân bố đều trên thang logarithm: Ridge trong khoảng $[1, 10^5]$ và Lasso trong khoảng $[10^{-5}, 10^1]$.
    \item Với mỗi ứng viên $\lambda$, lặp lại 5 lần: sử dụng 4 folds để huấn luyện và 1 fold còn lại làm tập kiểm thử cục bộ (validation set) để tính toán MSE.
    \item \textbf{Kỹ thuật Warm Start tối ưu hóa thời gian chạy:} Đối với thuật toán Coordinate Descent của Lasso, các ứng viên $\lambda$ được sắp xếp theo thứ tự giảm dần. Hệ số $\boldsymbol{\beta}$ tối ưu của bước chạy $\lambda_{t}$ trước đó được dùng làm điểm khởi tạo (warm start) cho bước chạy $\lambda_{t+1}$ tiếp theo, giúp giảm tới 80\% số vòng lặp cần thiết để Coordinate Descent hội tụ.
    \item Tính toán MSE trung bình và chọn $\lambda^*$ tối ưu có sai số nhỏ nhất để huấn luyện mô hình chính thức trên toàn bộ tập Train.
\end{enumerate}

\begin{figure}[H]
    \centering
    \begin{minipage}{0.48\textwidth}
        \centering
        \includegraphics[width=\textwidth]{../part2/figures/cv_curve_ridge.png}
        \caption{Đường cong Cross-Validation cho mô hình Ridge}
        \label{fig:cv_ridge}
    \end{minipage}\hfill
    \begin{minipage}{0.48\textwidth}
        \centering
        \includegraphics[width=\textwidth]{../part2/figures/cv_curve_lasso.png}
        \caption{Đường cong Cross-Validation cho mô hình Lasso}
        \label{fig:cv_lasso}
    \end{minipage}
\end{figure}

\subsection{Kết Quả Chạy Mô Hình Thực Tế}
Sau khi thực thi pipeline tiền xử lý tối ưu (rút gọn tập biến xuống còn \textbf{40 đặc trưng} thông qua gom nhóm định tính thời gian nhằm tối đa hóa tính giải thích và sự ổn định), cả 6 mô hình hồi quy đã được huấn luyện hoàn toàn từ đầu bằng NumPy (riêng Lasso được gọi qua sklearn để kiểm chứng do không có nghiệm đóng) và cho kết quả so sánh thực nghiệm trên tập Test vô cùng ấn tượng:

\begin{table}[H]
\centering
\caption{Bảng so sánh hiệu năng các mô hình trên tập Test (Beijing PM2.5)}
\rowcolors{2}{HCMUSBlue!5}{white}
\begin{tabular}{lccccc}
\toprule
\textbf{Mô hình hồi quy} & \textbf{MAE} & \textbf{RMSE} & \textbf{$R^2$ Test} & \textbf{Adj. $R^2$} & \textbf{Số đặc trưng} \\
\midrule
\textbf{Kernel Ridge (RBF Kernel)} & \textbf{0.2043} & \textbf{0.3238} & \textbf{0.9186} & \textbf{0.9179} & Phi tuyến \\
\textbf{OLS Basic} & \textbf{0.2043} & \textbf{0.3318} & \textbf{0.9146} & \textbf{0.9138} & 40 \\
\textbf{Bayesian Linear} & \textbf{0.2041} & \textbf{0.3319} & \textbf{0.9145} & \textbf{0.9138} & 40 \\
\textbf{Ridge ($\lambda = 16.77$)} & \textbf{0.2045} & \textbf{0.3319} & \textbf{0.9145} & \textbf{0.9137} & 40 \\
\textbf{Lasso ($\lambda = 0.0001265$)} & \textbf{0.2038} & \textbf{0.3320} & \textbf{0.9145} & \textbf{0.9137} & 40 (35 non-zero) \\
\textbf{OLS Feature Selected} & \textbf{0.2059} & \textbf{0.3354} & \textbf{0.9127} & \textbf{0.9119} & 25 \\
\bottomrule
\end{tabular}
\end{table}

\noindent\textbf{Nhận xét và phân tích chuyên sâu hiệu năng:}
\begin{enumerate}
    \item \textbf{Khả năng giải thích vượt trội của mô hình tuyến tính:} Tất cả các mô hình đều đạt chỉ số $R^2 > 0.91$ trên tập dữ liệu kiểm thử thực tế. Hơn 91.4\% sự biến thiên phức tạp của nồng độ bụi mịn PM2.5 đã được giải thích trọn vẹn bởi mô hình tuyến tính hồi quy kết hợp Feature Engineering mạnh mẽ.
    \item \textbf{Hiệu năng tuyệt vời của mô hình phi tuyến Kernel Ridge (RBF):} Kernel Ridge với RBF kernel đạt hiệu năng cao nhất ($R^2 = 91.86\%$ và $RMSE = 0.3238$). Nhờ phép chiếu hạt nhân phi tuyến tính lên không gian vô hạn chiều, mô hình bắt được các tương tác khí hậu vi mô cực kỳ mượt mà. Tuy nhiên, nó đòi hỏi chi phí tính toán rất lớn và không thể suy luận trực quan các hệ số hồi quy như mô hình OLS.
    \item \textbf{Tính tối ưu của OLS Feature Selected (loại bỏ ngược):} Bằng cách loại bỏ 15 biến kém quan trọng có $p$-value $\ge 0.05$ và chỉ giữ lại đúng \textbf{25 đặc trưng cốt lõi}, mô hình đạt hiệu năng rất cao với $R^2 \approx 91.27\%$ và $RMSE \approx 0.3354$. Việc giảm tới 37.5\% số lượng biến giúp mô hình cực kỳ tinh gọn, hạn chế hoàn toàn hiện tượng đa cộng tuyến và cực kỳ dễ giải thích mà sai số RMSE chỉ tăng một lượng vô cùng nhỏ ($\approx 0.0036$).
    \item \textbf{Vai trò lựa chọn đặc trưng tự động của Lasso:} Lasso với siêu tham số tối ưu lựa chọn qua Cross-Validation ($\lambda = 0.0001265$) đã triệt tiêu hệ số $\beta = 0$ của đúng 5 đặc trưng không quan trọng: \texttt{RAIN\_x\_WSPM}, \texttt{season\_Ha} (mùa Hè), \texttt{day\_of\_week\_1} (thứ Hai), \texttt{day\_of\_week\_2} (thứ Ba), và hướng gió \texttt{wd\_WSW}, giúp tinh giản mô hình một cách tự động mà vẫn giữ nguyên độ chính xác.
\end{enumerate}

\begin{figure}[H]
    \centering
    \includegraphics[width=0.9\textwidth]{../part2/figures/model_comparison.png}
    \caption{Biểu đồ so sánh hiệu năng các mô hình trên tập kiểm thử (MAE, RMSE, $R^2$)}
    \label{fig:model_comparison}
\end{figure}

\subsection{Đánh Giá Mô Hình \& Chẩn Đoán Phần Dư Thực Tế}

\subsubsection{Đường cong học tập và đánh giá overfitting}
\begin{figure}[H]
    \centering
    \includegraphics[width=0.75\textwidth]{../part2/figures/learning_curve.png}
    \caption{Đồ thị đường cong học tập của mô hình OLS}
\end{figure}

Đồ thị đường cong học tập thể hiện mối quan hệ giữa kích thước tập dữ liệu huấn luyện và điểm số R² / RMSE:
\begin{enumerate}
    \item \textbf{Tuyệt đối không có hiện tượng Overfitting:} Khoảng cách giữa hai đường Train và Test cực kỳ nhỏ hẹp và nhanh chóng hội tụ bám sát nhau khi lượng dữ liệu huấn luyện vượt quá mốc 30\%.
    \item \textbf{Năng lực học tập lý tưởng:} Cả hai đường đều duy trì ổn định ở mức rất cao ($> 0.91$), chứng tỏ mô hình có độ tổng quát hóa hoàn hảo trên dữ liệu tương lai chưa từng thấy và không có dấu hiệu bị Underfitting.
\end{enumerate}

\subsubsection{Chẩn đoán phần dư thực tế của mô hình tốt nhất}
Ta tiến hành vẽ 4 biểu đồ chẩn đoán phần dư thực tế của mô hình OLS để kiểm chứng các giả thiết Gauss-Markov:

\begin{figure}[H]
    \centering
    \begin{minipage}{0.48\textwidth}
        \centering
        \includegraphics[width=\textwidth]{../part2/figures/residuals_vs_fitted.png}
        \caption{Đồ thị Residuals vs Fitted}
    \end{minipage}\hfill
    \begin{minipage}{0.48\textwidth}
        \centering
        \includegraphics[width=\textwidth]{../part2/figures/qq_plot.png}
        \caption{Đồ thị Normal Q-Q Plot}
    \end{minipage}
\end{figure}

\begin{figure}[H]
    \centering
    \begin{minipage}{0.48\textwidth}
        \centering
        \includegraphics[width=\textwidth]{../part2/figures/scale_location.png}
        \caption{Đồ thị Scale-Location}
    \end{minipage}\hfill
    \begin{minipage}{0.48\textwidth}
        \centering
        \includegraphics[width=\textwidth]{../part2/figures/cooks_distance.png}
        \caption{Đồ thị Cook's Distance}
    \end{minipage}
\end{figure}

\noindent\textbf{Nhận xét từ các biểu đồ chẩn đoán:}
\begin{enumerate}
    \item \textbf{Residuals vs Fitted:} Đường trung bình màu đỏ nằm trùng khít lên trục hoành $y=0$, các điểm phần dư phân bổ đối xứng và ngẫu nhiên hoàn toàn. Điều này xác nhận giả thiết mối quan hệ tuyến tính (GM1) là hoàn toàn phù hợp và không bị bỏ sót dạng phi tuyến nào.
    \item \textbf{Normal Q-Q:} Đại đa số các điểm phần dư bám sát đường phân giác chuẩn tuyến tính, chỉ lệch nhẹ ở hai đầu cực trị. Điều này chứng minh giả thiết sai số phân phối chuẩn (GM5) được đáp ứng rất tốt, đảm bảo độ tin cậy tuyệt đối của các suy luận thống kê và khoảng tin cậy.
    \item \textbf{Scale-Location:} Đường đỏ trung bình phẳng lặng, các điểm phần dư phân tán đều không tạo thành hình phễu, xác nhận giả thiết phương sai sai số đồng đều không đổi (GM4) được thỏa mãn hoàn toàn.
    \item \textbf{Cook's Distance:} Nhận diện được một số điểm dị biệt có ảnh hưởng lớn vượt ngưỡng $4/n$ nhưng giá trị khoảng cách Cook tối đa vẫn ở mức an toàn ($D_i < 0.15$), đảm bảo không có bất kỳ điểm cá biệt nào thao túng làm lệch lạc hệ số hồi quy của toàn bộ hệ thống.
\end{enumerate}

\subsection{Tầm Quan Trọng của Đặc Trưng}
\begin{figure}[H]
    \centering
    \includegraphics[width=0.75\textwidth]{../part2/figures/feature_importance.png}
    \caption{Tầm quan trọng của các đặc trưng hồi quy (Top 20 Coefficients)}
\end{figure}

Dựa trên biểu đồ hệ số hồi quy chuẩn hóa (Feature Importance), ta đúc kết được các quy luật khí tượng học thực tế sâu sắc:
\begin{itemize}
    \item \textbf{Biến trễ tự hồi quy (\texttt{log\_PM2.5\_lag\_1}):} Có trọng số dương lớn nhất ($\beta \approx +0.8162$). Điều này phản ánh đặc trưng tự tương quan thời gian cực mạnh của khí quyển: nồng độ ô nhiễm của giờ trước đó là yếu tố tiên quyết quyết định nồng độ của giờ hiện tại.
    \item \textbf{Đặc trưng nhiệt động học (`DEWP` và `TEMP`):} Nhiệt độ điểm sương \texttt{DEWP} mang hệ số dương ($\beta \approx +0.0969$) do độ ẩm không khí cao tạo điều kiện cho các phản ứng hóa học hình thành bụi mịn thứ cấp. Nhiệt độ \texttt{TEMP} mang hệ số âm ($\beta \approx -0.0947$) thể hiện hiện tượng nghịch nhiệt vào mùa lạnh (không khí lạnh bị chìm xuống giữ bụi mịn ở tầng sát mặt đất).
    \item \textbf{Đặc trưng động học dòng khí (`WSPM`):} Tốc độ gió mang hệ số âm rất lớn ($\beta \approx -0.0897$), minh chứng cho cơ chế khuếch tán cơ học: gió càng mạnh thì bụi mịn càng bị pha loãng và phát tán nhanh chóng.
\end{itemize}

\clearpage

\section{Kết luận \& Đánh giá mức độ hoàn thành}

\subsection{Tóm tắt kết quả đồ án}
Đồ án đã hoàn thành xuất sắc toàn bộ các yêu cầu đặt ra từ cả hai khía cạnh lý thuyết học thuật và ứng dụng thực tế:
\begin{itemize}
    \item \textbf{Về mặt lý thuyết (Phần 1):} Hệ thống hóa và chứng minh toán học chặt chẽ các nền tảng của hồi quy tuyến tính OLS, các giả thiết Gauss-Markov và định lý BLUE. Thiết kế và cài đặt thành công thuật toán giải hệ Khử Gauss với Pivoting chọn trục xoay từ đầu bằng Python thuần (không dùng thư viện số học nâng cao cho thuật toán chính) để đảm bảo độ ổn định số học tối đa.
    \item \textbf{Về mặt thực tiễn (Phần 2):} Xây dựng một pipeline tiền xử lý dữ liệu chuẩn mực và giải quyết triệt để đa cộng tuyến hoàn hảo bằng phân rã QR. Huấn luyện thành công 6 mô hình hồi quy đa dạng với độ chính xác vô cùng ấn tượng ($R^2 > 0.91$), hoạt động cực kỳ ổn định không bị quá khớp và đáp ứng hoàn hảo các kiểm định chẩn đoán phần dư thực tế.
\end{itemize}

\subsection{Bài học rút ra}
Thông qua quá trình nghiên cứu và thực thi đồ án này, nhóm đã gặt hái được nhiều bài học học thuật quý báu:
\begin{itemize}
    \item \textbf{Kỹ năng lập trình toán học thuần chất:} Việc cài đặt các thuật toán OLS, Ridge Trace, giải hệ phương trình Khử Gauss từ đầu bằng Python thuần giúp nhóm thấu hiểu sâu sắc bản chất số học của máy tính, đặc biệt là các vấn đề tích tụ sai số làm tròn và tầm quan trọng của tính ổn định số học trong đại số tuyến tính thực hành.
    \item \textbf{Quy trình xử lý dữ liệu học thuật:} Khảo sát dữ liệu (EDA) chuyên sâu giúp thấu hiểu các phân phối thực tế của khí quyển. Việc xử lý khuyết thiếu bằng cách kết hợp giữa Interpolation, KNN Imputer và Mode Window chứng minh tầm quan trọng của việc lập luận khoa học dựa trên bản chất vật lý của dữ liệu thay vì áp dụng máy móc.
    \item \textbf{Tư duy mô hình hóa đa chiều:} Việc so sánh hiệu năng giữa OLS, Ridge, Lasso, Bayesian Linear và Kernel Ridge giúp nhóm phân biệt rõ ràng giữa phương pháp tuyến tính và phi tuyến, cơ chế kiểm soát quá khớp (overfitting) bằng các kỹ thuật điều chuẩn (regularization).
\end{itemize}

\subsection{Hướng mở rộng phát triển}
Dựa trên các kết quả đạt được, mô hình dự báo chất lượng không khí có thể được mở rộng theo các hướng nghiên cứu sau:
\begin{itemize}
    \item \textbf{Tích hợp cấu trúc thời gian phức tạp:} Áp dụng các mô hình chuỗi thời gian chuyên sâu như SARIMAX hoặc các mạng nơ-ron hồi quy recurrent (LSTM, GRU) để khai thác sâu hơn mối tương quan trễ và các chu kỳ dài hạn trong khí quyển.
    \item \textbf{Mở rộng không gian (Multi-site Modeling):} Hiện tại mô hình chỉ dự báo cho trạm Shunyi. Hướng tiếp theo có thể xây dựng mô hình không-thời gian (spatio-temporal models) kết hợp dữ liệu từ tất cả 12 trạm khí tượng tại Bắc Kinh nhằm nắm bắt quy luật lan truyền ô nhiễm theo hướng gió giữa các vùng địa lý.
    \item \textbf{Khai phá các đặc trưng nâng cao:} Đưa thêm các thông số về mật độ giao thông, hoạt động công nghiệp, bản đồ nhiệt độ đô thị để nâng cao độ chính xác dự báo trong các điều kiện thời tiết cực đoan.
\end{itemize}

\subsection{Tự đánh giá mức độ hoàn thành}
Dưới đây là các bảng tự đánh giá chi tiết mức độ hoàn thành đối với từng phần yêu cầu được đặt ra trong đề bài Đồ án 2.

\begin{table}[H]
    \centering
    \caption{Bảng đánh giá mức độ hoàn thành các yêu cầu Phần 1 (Lý thuyết và Minh họa)}
    \rowcolors{2}{HCMUSBlue!5}{white}
    \resizebox{\textwidth}{!}{
    \begin{tabular}{cp{3.5cm}p{7cm}cc}
        \toprule
        \textbf{Phần} & \textbf{Tiêu chí} & \textbf{Minh chứng chi tiết từ Mã nguồn \& Báo cáo} & \textbf{Điểm} & \textbf{Hoàn thành} \\
        \midrule
        \textbf{1} & Lý thuyết OLS & Đầy đủ công thức, chứng minh Normal Equations bằng đạo hàm & 1.0 & 100\% \\
        \textbf{1} & Cài đặt OLS & Hàm \texttt{ols\_fit(X, y)} từ đầu dùng khử Gauss có pivoting & 1.0 & 100\% \\
        \textbf{1} & Hat matrix & Hàm \texttt{hat\_matrix(X)}, kiểm chứng tính chất đối xứng \& idempotent & 0.5 & 100\% \\
        \textbf{1} & Kiểm định hệ số & Thống kê $t$, $F$, trị số p-value chuẩn xác (\texttt{coef\_inference}) & 0.5 & 100\% \\
        \textbf{1} & Regularization & Hồi quy Ridge (\texttt{ridge\_fit}) và hàm vẽ Ridge Trace (\texttt{ridge\_trace}) & 1.0 & 100\% \\
        \textbf{1} & Phân tích phần dư & Vẽ đủ 4 biểu đồ chẩn đoán chuẩn mực trong \texttt{residual\_analysis.py} & 0.5 & 100\% \\
        \textbf{1} & Cross-validation & Tự cài đặt hàm k-Fold CV (\texttt{kfold\_cv}) không sử dụng thư viện ngoài & 0.5 & 100\% \\
        \textbf{1} & Minh họa BLUE & Mô phỏng Monte Carlo 1000 lần chứng minh tính không chệch & 0.5 & 100\% \\
        \textbf{1} & Trình bày Notebook & Notebook \texttt{part1\_notebook.ipynb} chi tiết, giải thích rõ ràng & 0.5 & 100\% \\
        \bottomrule
    \end{tabular}
    }
\end{table}

\begin{table}[H]
    \centering
    \caption{Bảng đánh giá mức độ hoàn thành các yêu cầu Phần 2 (Ứng dụng thực tế)}
    \rowcolors{2}{HCMUSBlue!5}{white}
    \resizebox{\textwidth}{!}{
    \begin{tabular}{cp{3.5cm}p{7cm}cc}
        \toprule
        \textbf{Phần} & \textbf{Tiêu chí} & \textbf{Minh chứng chi tiết từ Mã nguồn \& Báo cáo} & \textbf{Điểm} & \textbf{Hoàn thành} \\
        \midrule
        \textbf{2} & Chọn bộ dữ liệu & Bộ dữ liệu PRSA trạm Shunyi ($n = 35.064, p = 11$), thỏa tiêu chí & 0.5 & 100\% \\
        \textbf{2} & Khảo sát EDA & Thống kê mô tả đầy đủ, vẽ histogram, boxplot, heatmap chi tiết & 0.5 & 100\% \\
        \textbf{2} & Xử lý missing & Giải thích cơ chế MCAR/MAR; áp dụng Linear Interp, KNN, Mode & 1.0 & 100\% \\
        \textbf{2} & Tiền xử lý & Class \texttt{DataPipeline} đóng gói \texttt{fit}/\texttt{transform} chống rò rỉ dữ liệu & 0.5 & 100\% \\
        \textbf{2} & Xây dựng mô hình & Huấn luyện và so sánh OLS, Lasso, Ridge, Polynomial, Bayesian & 1.5 & 100\% \\
        \textbf{2} & Đánh giá mô hình & Đánh giá bằng MAE, RMSE, $R^2$ trên Test; chẩn đoán 4 biểu đồ phần dư & 1.0 & 100\% \\
        \textbf{2} & Nhận xét kết luận & Giải thích chi tiết các kết quả dựa trên đặc thù môi trường Bắc Kinh & 0.5 & 100\% \\
        \textbf{2} & Kỹ thuật nâng cao & Bayesian Linear \& Kernel Ridge from scratch đạt \textbf{+0.5 điểm thưởng} & +0.5 & \textbf{Xuất sắc} \\
        \bottomrule
    \end{tabular}
    }
\end{table}

\clearpage

\phantomsection
\addcontentsline{toc}{section}{Tài liệu tham khảo}
\bibliographystyle{plain}
\bibliography{report_final}

\end{document}
